\input{/Users/tommasopedroni/Desktop/TESI/Tesi/latex_prima_versione/preambolo_bozza.tex}
\newgeometry{margin=2cm}
\setlength{\headwidth}{\textwidth}

\begin{document}
\setcounter{chapter}{3}
\chapter{Von Neumann Algebras and The Modular Operator}
\section{Von Neumann Algebras, Concretely and Abstractly}\label{sec:vn-algebras}

The previous chapter derived Schr\"odinger dynamics, the Born rule, and
Noether's theorem from the intrinsic K\"ahler geometry of the projective
Hilbert space, once a state and a distinguished selfadjoint Hamiltonian
were given. The role of the Hamiltonian itself, however, was never
questioned. In every construction so far, a time evolution enters the
theory as external structure, whether as the one-parameter group of
Definition 1.6.23 in the classical setting or as the Hamiltonian flow of
Proposition 3.3.2 in the quantum one. This chapter asks whether that
structure can instead be recovered from the algebra together with a
single state, without postulating a Hamiltonian at all. The
Tomita--Takesaki theorem shows that every faithful normal state on a von
Neumann algebra determines a canonical one-parameter group of
automorphisms, its modular automorphism group. When the state describes
thermal equilibrium, this group coincides with ordinary time evolution
up to a rescaling by the inverse temperature, an identification that
suggests reading the modular flow itself as the origin of physical time
whenever no external clock is available. This is the content of the
thermal time hypothesis of Connes and Rovelli, and it is the
algebraic-geometric standpoint developed throughout this thesis that
makes the proposal precise.

Every construction so far in this thesis singled out the $C^*$-algebra
of Definition 1.6.11 as the primitive algebraic object, and derived a
Hilbert space representation of an observable algebra only afterward,
through the Gelfand-Naimark-Segal construction of Theorem 2.7. The
states considered there were arbitrary positive normalized functionals,
with no continuity required beyond what automatic continuity of
Proposition 2.1.23 already supplies. Tomita--Takesaki theory, together
with the notion of a thermal state occupying Section 4.5, requires a
strictly stronger setting. The relevant algebra of observables must be
weakly closed, so that it coincides with its own bicommutant, and the
state that generates the modular flow must be normal, so that
expectation values remain continuous under the monotone limits used to
define equilibrium. This section introduces von Neumann algebras and
normal states as this strengthened setting, and isolates the property
under which the Gelfand-Naimark-Segal vector of Theorem 2.7 is
simultaneously cyclic for the algebra and separating for it, a
distinction the construction of Section 4.3 will use directly. A
concrete algebra of operators on a fixed Hilbert space, however, is not
the natural datum on the abstract side of the correspondence developed
in Chapters 1 and 2, where the observable algebra $A$ of Definition
2.1.22 carries no reference to any particular Hilbert space. The second
half of this section therefore isolates the property of $A$ alone,
having a predual, that identifies exactly when $A$ can be realized this
way, and closes a gap left open by the Gelfand-Naimark-Segal
construction itself: nothing in Theorem 2.7 guaranteed that the
represented algebra $\pi_\omega(A)$ is weakly closed, hence a genuine
von Neumann algebra in the sense about to be introduced.

\begin{definition}[Commutant]\label{def:commutant}
	Let $H$ be a Hilbert space, as per Definition 2.2.6, and let
	$S \subseteq B(H)$ be a subset. The commutant of $S$ is
	\[
	S' := \{T \in B(H) \mid TA = AT \ \forall A \in S\}.
	\]
\end{definition}

\begin{remark}
	The commutant $S'$ is a subalgebra of $B(H)$ containing the identity,
	and every element of $S$ commutes with every element of $S'$ by
	definition, so $S$ is contained in its double commutant $S'' := (S')'$.
\end{remark}

The characterization of an algebra through its double commutant, rather
than through an explicit topological closure, is only useful once the
relevant topology on $B(H)$ has itself been fixed. The operator norm
topology used throughout this thesis is too fine for this purpose,
since it does not detect the monotone limits under which a state is
required to behave continuously in Definition~\ref{def:normal-state}
below, so a coarser pair of topologies is introduced first.

\begin{definition}[Weak and strong operator topologies]\label{def:wot-sot}
	Let $H$ be a Hilbert space, as per Definition 2.2.6. The weak operator
	topology on $B(H)$ is the coarsest topology for which every map
	$T \mapsto (\phi, T\psi)$, with $\phi, \psi \in H$, is continuous. The
	strong operator topology is the coarsest topology for which every map
	$T \mapsto \lVert T\psi \rVert$, with $\psi \in H$, is continuous.
\end{definition}

With both topologies available, it becomes possible to ask whether the
algebraic condition $M = M''$ used above coincides with topological
closure. Von Neumann's original theorem answers this affirmatively for
unital $*$-subalgebras, and this equivalence is what justifies treating
the double commutant, an entirely algebraic condition, as the working
definition of the topologically closed algebras this chapter needs.

\begin{theorem}[von Neumann Bicommutant Theorem]\label{thm:bicommutant}
	Let $M \subseteq B(H)$ be a unital $*$-subalgebra, as per
	Definition 1.6.2. The following are equivalent.
	\begin{enumerate}
		\item $M = M''$, where $M''$ denotes the double commutant of
		Definition~\ref{def:commutant}.
		\item $M$ is closed in the weak operator topology of
		Definition~\ref{def:wot-sot}.
		\item $M$ is closed in the strong operator topology of
		Definition~\ref{def:wot-sot}.
	\end{enumerate}
\end{theorem}
\noindent This is proved in \cite[TODO: collocazione precisa]{TODO-vonNeumannBicommutant}.

\begin{definition}[von Neumann algebra]\label{def:vN-algebra}
	A von Neumann algebra on $H$ is a unital $*$-subalgebra
	$M \subseteq B(H)$, as per Definition 1.6.2, satisfying $M = M''$, as
	per Definition~\ref{def:commutant}.
\end{definition}

\begin{remark}
	By Theorem~\ref{thm:bicommutant}, a von Neumann algebra can equally be
	characterized as a unital $*$-subalgebra of $B(H)$ closed in either
	operator topology of Definition~\ref{def:wot-sot}. Every von Neumann
	algebra is in particular a $C^*$-algebra, as per Definition 1.6.11,
	under the operator norm, since the norm topology is finer than both
	operator topologies and closure therefore persists from either of them
	to the norm topology.
\end{remark}

The states relevant to modular theory must interact well with the
operator topologies of Definition~\ref{def:wot-sot}, not merely with
the norm topology already available on any $C^*$-algebra by Definition
2.1.3. The most direct way to express this compatibility is to exhibit
each such state as the restriction of a single linear functional
defined on the whole of $B(H)$, built from a fixed trace-class operator
rather than from the algebra $M$ itself. This requires first extending
the trace of Definition 2.3.15, so far used only for finite-rank
density matrices in Proposition 2.3.14, to a linear functional defined
on a genuine two-sided ideal of $B(H)$.

\begin{definition}[Trace of a positive operator]\label{def:trace}
	Let $H$ be a Hilbert space, as per Definition 2.2.6, and let
	$\{e_i\}_{i \in I}$ be an orthonormal basis of $H$. For $A \in B(H)$
	positive, as per Definition 2.2.22, the trace of $A$ is
	\[
	\mathrm{Tr}(A) := \sum_{i \in I} (e_i, A e_i) \in [0, \infty].
	\]
\end{definition}

\begin{remark}
	The value of $\mathrm{Tr}(A)$ in Definition~\ref{def:trace} does not
	depend on the choice of orthonormal basis $\{e_i\}_{i\in I}$, so that
	$\mathrm{Tr}$ is a well defined map from the positive elements of
	$B(H)$ to $[0,\infty]$. This invariance is proved in \cite[TODO:
	collocazione precisa]{TODO-traceclass}.
\end{remark}

Definition~\ref{def:trace} assigns a value in $[0,\infty]$ to every
positive operator, but a genuine linear functional requires this value
to be finite everywhere, and a genuine ideal of $B(H)$ requires closure
not only under sums of positive elements but under multiplication by an
arbitrary bounded operator on either side. Both properties hold once
the trace is imposed on the modulus of a general, not necessarily
positive, operator.

\begin{definition}[Trace-class operators]\label{def:trace-class}
	Let $H$ be a Hilbert space, as per Definition 2.2.6, and let
	$A \in B(H)$. Write $|A| := (A^*A)^{1/2}$, the positive square root of
	$A^*A$ as per Remark 2.2.23, positive since
	$(\psi, A^*A\psi) = \lVert A\psi \rVert^2 \geq 0$ for every $\psi \in
	H$. The operator $A$ is trace-class if $\mathrm{Tr}(|A|) < \infty$, as
	per Definition~\ref{def:trace}. The set of trace-class operators on $H$
	is denoted $B_1(H)$.
\end{definition}

\begin{remark}
	The notation $B_1(H)$ parallels the notation $B_0(H)$ of Definition
	2.2.15 for the ideal of compact operators. In fact, $B_1(H)$ is itself
	a two-sided ideal of $B(H)$, contained in $B_0(H)$, and the trace
	extends from the positive elements of $B_1(H)$ to a linear functional
	on the whole of $B_1(H)$, satisfying
	\[
	\mathrm{Tr}(AB) = \mathrm{Tr}(BA), \qquad A \in B_1(H), \ B \in B(H).
	\]
	Setting $\lVert A \rVert_1 := \mathrm{Tr}(|A|)$ makes
	$(B_1(H), \lVert \cdot \rVert_1)$ a Banach space, as per Definition
	1.6.9, and the map
	\[
	B(H) \to (B_1(H))', \qquad T \mapsto \big( S \mapsto \mathrm{Tr}(ST)
	\big),
	\]
	with $(B_1(H))'$ the topological dual of Definition 2.2.13, is an
	isometric isomorphism.\label{eq:trace-duality-remark} Every fact
	recalled in this remark is proved in \cite[TODO: collocazione
	precisa]{TODO-traceclass}.
\end{remark}

The isometric identification of $B(H)$ with the topological dual of
$B_1(H)$ singles out a topology on $B(H)$ coarser than the operator
norm, and, as the following remark records, also finer than the two
topologies of Definition~\ref{def:wot-sot} in general, yet coincident
with them on the bounded sets relevant to representing states.

\begin{definition}[Ultraweak topology]\label{def:ultraweak}
	Let $H$ be a Hilbert space, as per Definition 2.2.6. The ultraweak
	topology on $B(H)$ is the coarsest topology for which every evaluation
	map $\widehat{S} : B(H) \to \mathbb{C}$, $\widehat{S}(T) :=
	\mathrm{Tr}(ST)$, with $S \in B_1(H)$ as per Definition
	\ref{def:trace-class}, is continuous, generalizing to the Banach space
	$B_1(H)$ the same weak$^*$ construction that Definition 1.6.28 performs
	on the topological dual of a $C^*$-algebra.
\end{definition}

\begin{remark}
	The ultraweak topology of Definition~\ref{def:ultraweak} is finer than
	the weak operator topology of Definition~\ref{def:wot-sot}, since the
	latter only tests continuity against rank-one operators $S =
	(\cdot,\phi)\psi$, a proper subset of $B_1(H)$. The two topologies
	coincide on norm-bounded subsets of $B(H)$, so that a von Neumann
	algebra $M$, as per Definition~\ref{def:vN-algebra}, closed in one is
	automatically closed in the other, consistently with
	Theorem~\ref{thm:bicommutant} already identifying both with the same
	purely algebraic condition. This equivalence is proved in \cite[TODO:
	collocazione precisa]{TODO-traceclass}.
\end{remark}

Restricting the ultraweakly continuous functionals of
Definition~\ref{def:ultraweak} from the whole of $B(H)$ to a von
Neumann subalgebra $M$ produces a distinguished Banach space attached
to $M$ itself, one that supplies the correct, structural, notion of a
normal state below.

\begin{definition}[Predual]\label{def:predual}
	Let $M$ be a von Neumann algebra on $H$, as per
	Definition~\ref{def:vN-algebra}. The predual of $M$ is the space
	\[
	M_* := \big\{ \omega|_M \mid \omega \in B_1(H) \big\}
	\]
	of restrictions to $M$ of the trace functionals $S \mapsto
	\mathrm{Tr}(S \, \cdot \,)$, as per Definition~\ref{def:trace-class},
	equipped with the quotient norm inherited from $B_1(H)$.
\end{definition}

\begin{remark}
	The predual $M_*$ of Definition~\ref{def:predual} is a Banach space, as
	per Definition 1.6.9, and the pairing $(\phi, A) \mapsto \phi(A)$, for
	$\phi \in M_*$ and $A \in M$, identifies $M$ isometrically with the
	topological dual $(M_*)'$ of $M_*$, as per Definition 2.2.13, and this
	identification is unique. Both facts are proved in \cite[TODO:
	collocazione precisa]{TODO-predual}.
\end{remark}

With the predual in hand, normality of a state can be expressed
intrinsically, as membership in $M_*$, rather than through the monotone
limits that only characterize it after the fact.

\begin{definition}[Normal state]\label{def:normal-state}
	Let $M$ be a von Neumann algebra on $H$, as per
	Definition~\ref{def:vN-algebra}. A state $\omega$ on $M$, as per
	Definition 2.1.3, is normal if $\omega \in M_*$, the predual of $M$ as
	per Definition~\ref{def:predual}.
\end{definition}

\begin{remark}
	Definition~\ref{def:normal-state} is equivalent to the more classical
	requirement that
	\[
	\omega\Big(\sup_\alpha A_\alpha\Big) = \sup_\alpha \omega(A_\alpha),
	\]
	for every bounded increasing net $(A_\alpha)$ of positive elements of
	$M$, with the supremum on the left taken in the strong operator
	topology of Definition~\ref{def:wot-sot}, and equally to continuity of
	$\omega$ for the ultraweak topology of Definition~\ref{def:ultraweak}
	restricted to $M$. This equivalence is proved in \cite[TODO:
	collocazione precisa]{TODO-normalStates}. When $M = B(H)$, it recovers
	the density-matrix description of Definition 2.3.15: normality of
	$\omega$ is equivalent to the existence of $\rho \in B_1(H)$ positive
	with $\mathrm{Tr}(\rho) = 1$ such that $\omega(A) = \mathrm{Tr}(\rho A)$
	for every $A \in B(H)$, matching Proposition 2.3.14.
\end{remark}

Every von Neumann algebra of Definition~\ref{def:vN-algebra} carries a
predual by construction, since Definition~\ref{def:predual} builds
$M_*$ directly from the ambient Hilbert space $H$. The converse
question, whether an abstract $C^*$-algebra with no reference to any
Hilbert space can already be recognized as a von Neumann algebra from
the existence of some predual alone, is the content of the abstract
characterization due to Sakai, and it is the form in which the
correspondence with Chapters 1 and 2 is used from Section 4.2 onward.

\begin{definition}[$W^*$-algebra]\label{def:wstar-algebra}
	A $W^*$-algebra is a $C^*$-algebra $A$, as per Definition 1.6.11, for
	which there exists a Banach space $A_*$, as per Definition 1.6.9, such
	that $A$ is isometrically isomorphic, as a Banach space, to the
	topological dual $(A_*)'$ of $A_*$, as per Definition 2.2.13.
\end{definition}

\begin{theorem}[Sakai]\label{thm:sakai}
	Let $A$ be a $W^*$-algebra, as per Definition~\ref{def:wstar-algebra}.
	There exist a Hilbert space $H$ and a $*$-isomorphism, as per
	Definition 2.3.1, of $A$ onto a von Neumann algebra $M$ on $H$, as per
	Definition~\ref{def:vN-algebra}. Moreover, the Banach space $A_*$ of
	Definition~\ref{def:wstar-algebra} is then unique, and this isomorphism
	identifies it with the predual $M_*$ of Definition~\ref{def:predual}.
\end{theorem}
\noindent This is proved in \cite[TODO: collocazione precisa]{TODO-Sakai}.

Theorem~\ref{thm:sakai} addresses the abstract side of the
correspondence, but it leaves open a question specific to the
construction already carried out in Chapter 2: nothing in the
Gelfand-Naimark-Segal construction of Theorem 2.7 was shown to produce
a weakly closed algebra of operators, so $\pi_\omega(A)$ was, until now,
only a $C^*$-subalgebra of $B(H_\omega)$, not yet known to be a von
Neumann algebra in the sense of Definition~\ref{def:vN-algebra}.
Normality of the generating state closes exactly this gap.

\begin{proposition}[GNS representations of normal states are weakly
	closed]\label{prop:gns-normal}
	Let $M$ be a von Neumann algebra on $H$, as per
	Definition~\ref{def:vN-algebra}, and let $\omega$ be a normal state on
	$M$, as per Definition~\ref{def:normal-state}. Let $(H_\omega,
	\pi_\omega, \Psi_\omega)$ be the Gelfand-Naimark-Segal representation
	of $\omega$, as per Theorem 2.7. Then $\pi_\omega(M)$ is a von Neumann
	algebra on $H_\omega$, as per Definition~\ref{def:vN-algebra}.
\end{proposition}
\noindent This is proved in \cite[TODO: collocazione precisa]{TODO-GNSnormal}.

\begin{remark}
	Proposition~\ref{prop:gns-normal} is what allows Sections 4.2 onward to
	treat $\pi_\omega(M)$ itself as a concrete von Neumann algebra, with
	$\Psi_\omega$ cyclic for it by part i) of Theorem 2.7, rather than
	merely as a $C^*$-subalgebra of $B(H_\omega)$ to which
	Theorem~\ref{thm:bicommutant} would not yet apply.
\end{remark}

The vector $\Psi_\omega$ obtained this way is cyclic by construction,
and its associated vector state is normal by the following general
fact, which holds for an arbitrary vector in a von Neumann algebra
without any reference to the Gelfand-Naimark-Segal construction itself.

\begin{remark}[Vector states are normal]\label{rem:vector-normal}
	Let $M$ be a von Neumann algebra on $H$, as per
	Definition~\ref{def:vN-algebra}, and let $\Psi \in H$ be a unit vector.
	Setting $\rho := P_\Psi$, the rank-one projector of Definition 2.2.18,
	the vector state $\omega(A) := (\Psi, A\Psi)$ coincides with the
	restriction to $M$ of the trace functional $A \mapsto \mathrm{Tr}(\rho
	A)$, as per Definition~\ref{def:trace-class}, since $\mathrm{Tr}(\rho
	A) = (\Psi, A\Psi)$ by the same computation as in Proposition 2.3.14.
	Hence every vector state on $M$ lies in the predual $M_*$ of
	Definition~\ref{def:predual}, and is in particular normal.
\end{remark}

Separation for $\pi_\omega(M)$ is a strictly stronger requirement than
cyclicity, and it is exactly the combination of both properties in a
single vector, on a von Neumann algebra rather than a bare
$C^*$-algebra, that Tomita's construction in Section 4.3 takes as its
point of departure.

\begin{lemma}[Cyclic and separating vectors]\label{lem:cyclic-separating}
	Let $M$ be a von Neumann algebra on $H$, as per
	Definition~\ref{def:vN-algebra}, and let $\Omega \in H$ be a unit
	vector, cyclic for $M$ as per Definition 2.2.34. Let $\omega$ be the
	associated vector state, $\omega(A) := (\Omega, A\Omega)$ for
	$A \in M$. Then $\Omega$ is separating for $M$, meaning that
	$A\Omega = 0$ with $A \in M$ forces $A = 0$, if and only if $\omega$ is
	faithful, as per Definition 2.3.17.
\end{lemma}

\begin{proof}
	Suppose $\omega$ is faithful and let $A \in M$ satisfy $A\Omega = 0$.
	Then $\omega(A^*A) = (\Omega, A^*A\Omega) = \lVert A\Omega \rVert^2 =
	0$, and faithfulness forces $A^*A = 0$. The $C^*$-identity of
	Definition 1.6.11 then yields $\lVert A \rVert^2 = \lVert A^*A \rVert =
	0$, that is, $A = 0$. Conversely, suppose $\omega$ is not faithful.
	There exists a nonzero $B \in M$ with $\omega(B^*B) = 0$, and the same
	computation gives $\lVert B\Omega \rVert^2 = \omega(B^*B) = 0$, so
	$B\Omega = 0$ with $B \neq 0$. Hence $\Omega$ fails to be separating
	for $M$.
\end{proof}

Remark 2.3.18 already showed that the Gelfand-Naimark-Segal
representation of a faithful state is itself faithful, by tracking the
kernel of the representation directly.
Lemma~\ref{lem:cyclic-separating} isolates the stronger statement
needed here, that the algebra acts on its own cyclic vector without
collapsing any nonzero element to zero, since it is this separating
property, rather than injectivity of the representation alone, that
Tomita's construction in Section 4.3 takes as its starting hypothesis.

\section{Factors, Projections, and the Trace Problem}\label{sec:factors}

Section 4.1 fixed the ambient algebraic setting for the rest of this
chapter, a von Neumann algebra $M$ together with a cyclic vector, and
the notion of a normal state guaranteeing that expectation values on
$M$ interact correctly with the operator topologies of
Definition~\ref{def:wot-sot}. Before constructing the modular flow
itself, it is worth understanding why a construction already available
in this thesis, representing a state through a density matrix as in
Definition 2.3.15 and Proposition 2.3.14, cannot serve as the general
framework this chapter needs. That construction rests on the trace of
Definition~\ref{def:trace} being available on the whole of $B(H)$, and
on the possibility of writing every normal state as $\omega(A) =
\mathrm{Tr}(\rho A)$ for some positive $\rho$ with $\mathrm{Tr}(\rho) =
1$. This section shows that such a trace need not exist on an
arbitrary von Neumann algebra at all, and classifies exactly when it
does, through the internal structure of the lattice of projections of
$M$. The classification singles out a family of algebras, of type III,
on which no nonzero trace of any kind is available, so that density
matrices in the sense of Chapter 2 simply do not exist there. It is
precisely this obstruction, rather than a technical inconvenience, that
motivates replacing the trace with the modular automorphism group
constructed in the remainder of the chapter, a substitute that requires
only a cyclic and separating vector, as isolated in
Lemma~\ref{lem:cyclic-separating}, and that reduces to the usual Gibbs
description of Section 4.6 whenever a trace happens to be available.

The comparison of projections that drives this classification is
carried out through operators that behave isometrically on part of $H$
and vanish on the rest, generalizing the unitary operators already
introduced in Definition 2.2.22 to the case where the initial and final
spaces need not coincide.

\begin{definition}[Partial isometry]\label{def:partial-isometry}
	Let $H$ be a Hilbert space, as per Definition 2.2.6. An operator
	$U \in B(H)$ is a partial isometry if $U^*U$ is a projector, as per
	Definition 2.2.17.
\end{definition}

\begin{remark}[Initial and final projections]\label{rem:partial-isometry}
	Let $U \in B(H)$ be a partial isometry, as per
	Definition~\ref{def:partial-isometry}, and write $P := U^*U$, the
	initial projection of $U$. For every $\psi \in H$,
	\[
	\lVert U\psi \rVert^2 = (\psi, U^*U\psi) = (\psi, P\psi) = \lVert
	P\psi \rVert^2,
	\]
	using $P^* = P = P^2$ as per Definition 2.2.17. Hence $\ker U = \ker
	P$, so that $U$ vanishes on $\ker P = (\operatorname{ran} P)^\perp$
	and restricts to an isometry on $\operatorname{ran}(P)$, a closed
	subspace by Theorem 2.4. Since $U$ vanishes on $\ker P$, it satisfies
	$UP = U$ as operators on $H$, so that $Q := UU^*$ satisfies
	\[
	Q^2 = UU^*UU^* = U(U^*U)U^* = UPU^* = UU^* = Q,
	\]
	while $Q^* = Q$ holds trivially. Hence $Q$ is itself a projector, as
	per Definition 2.2.17, called the final projection of $U$, with range
	equal to $\operatorname{ran}(U)$.
\end{remark}

Partial isometries supply exactly the notion of comparison needed to
declare two projections of a von Neumann algebra equal in size, even
when they are not literally the same operator, a comparison that
underlies the whole classification to follow.

\begin{definition}[Murray-von Neumann equivalence]\label{def:mv-equivalence}
	Let $M$ be a von Neumann algebra on $H$, as per
	Definition~\ref{def:vN-algebra}. Two projectors $P, Q \in M$, as per
	Definition 2.2.17, are Murray-von Neumann equivalent, written $P \sim
	Q$, if there exists a partial isometry $U \in M$, as per
	Definition~\ref{def:partial-isometry}, with $U^*U = P$ and $UU^* = Q$.
\end{definition}

\begin{remark}
	Definition~\ref{def:mv-equivalence} defines an equivalence relation on
	the projectors of $M$. Reflexivity holds since $P$ itself is a partial
	isometry with $P^*P = P^2 = P$, so $U = P$ witnesses $P \sim P$.
	Symmetry holds since $U$ witnessing $P \sim Q$ forces $U^*$, with
	$(U^*)^*U^* = UU^* = Q$ and $U^*(U^*)^* = U^*U = P$, to witness $Q \sim
	P$. For transitivity, let $U_1$ witness $P \sim Q$ and $U_2$ witness $Q
	\sim R$. Since $Q$ is the final projection of $U_1$ and the initial
	projection of $U_2$, Remark~\ref{rem:partial-isometry} yields $QU_1 =
	U_1$ and $U_2Q = U_2$, so that $U := U_2U_1 \in M$ satisfies
	\[
	U^*U = U_1^*U_2^*U_2U_1 = U_1^*QU_1 = (QU_1)^*U_1 = U_1^*U_1 = P,
	\qquad UU^* = U_2U_1U_1^*U_2^* = U_2QU_2^* = U_2U_2^* = R,
	\]
	so $U$ witnesses $P \sim R$.
\end{remark}

Murray-von Neumann equivalence compares projections without reference
to any external notion of size, but the classification of factors also
needs to compare a projection with a strictly smaller one contained
inside it, a containment already implicit in the range inclusions used
above.

\begin{remark}[Order on projectors]\label{rem:projector-order}
	For projectors $P, Q \in M$, as per Definition 2.2.17, write $P \leq
	Q$ when $PQ = P$. This holds if and only if $\operatorname{ran}(P)
	\subseteq \operatorname{ran}(Q)$, and it turns the set of projectors of
	$M$ into a partially ordered set.
\end{remark}

With both a notion of equivalence and a notion of containment
available, a projection can be tested against its own equivalents
lying strictly inside it, and the outcome of this test is what
separates the finite from the infinite case in the trichotomy below.

\begin{definition}[Finite and infinite projections]\label{def:finite-projection}
	Let $M$ be a von Neumann algebra on $H$, as per
	Definition~\ref{def:vN-algebra}, and let $P \in M$ be a projector, as
	per Definition 2.2.17. The projector $P$ is finite if every projector
	$Q \in M$ satisfying $Q \leq P$, as per Remark~\ref{rem:projector-order},
	and $Q \sim P$, as per Definition~\ref{def:mv-equivalence}, satisfies
	$Q = P$. The projector $P$ is infinite if it is not finite.
\end{definition}

Having compared projections against each other, the classification
turns to the algebra as a whole, through the elements that commute
with every one of its own members.

\begin{definition}[Center]\label{def:center}
	Let $M$ be a von Neumann algebra on $H$, as per
	Definition~\ref{def:vN-algebra}. The center of $M$ is $Z(M) := M \cap
	M'$, with $M'$ the commutant of $M$, as per Definition~\ref{def:commutant}.
\end{definition}

\begin{remark}
	The center $Z(M)$ of Definition~\ref{def:center} is itself abelian.
	Indeed, for $A, B \in Z(M)$, $A \in M$ and $B \in Z(M) \subseteq M'$
	force $AB = BA$, since every element of $M'$ commutes with every
	element of $M$ by Definition~\ref{def:commutant}.
\end{remark}

\begin{definition}[Factor]\label{def:factor}
	A von Neumann algebra $M$ on $H$, as per Definition~\ref{def:vN-algebra},
	is a factor if $Z(M) = \mathbb{C}1$, with $Z(M)$ the center of $M$, as
	per Definition~\ref{def:center}.
\end{definition}

The factors are the von Neumann algebras with no nontrivial invariant
splitting into independent pieces, and every von Neumann algebra
decomposes into a direct integral of factors. The trichotomy below
concerns exactly these building blocks. Its statement, however,
requires generalizing the finite-dimensional trace already used
informally in Definition 2.3.15 to an object that need not be bounded
and need not even be finite, so that it can register the difference
between a factor on which such a functional exists and one on which it
does not.

\begin{definition}[Trace]\label{def:trace-functional}
	Let $M$ be a von Neumann algebra on $H$, as per
	Definition~\ref{def:vN-algebra}, and let $M_+$ denote its positive
	elements, as per Definition 2.2.22. A trace on $M$ is a map $\tau :
	M_+ \to [0, \infty]$ satisfying
	\[
	\tau(A + B) = \tau(A) + \tau(B), \qquad \tau(\lambda A) = \lambda
	\tau(A), \qquad \tau(UAU^*) = \tau(A),
	\]
	for all $A, B \in M_+$, all $\lambda \geq 0$, with the convention $0
	\cdot \infty := 0$, and all unitary $U \in M$, as per Definition 2.2.22.
\end{definition}

\begin{definition}[Faithful trace]\label{def:faithful-trace}
	Let $\tau$ be a trace on $M$, as per Definition~\ref{def:trace-functional}.
	The trace $\tau$ is faithful if $\tau(A) = 0$ with $A \in M_+$ forces
	$A = 0$.
\end{definition}

\begin{definition}[Normal trace]\label{def:normal-trace}
	Let $\tau$ be a trace on $M$, as per Definition~\ref{def:trace-functional}.
	The trace $\tau$ is normal if
	\[
	\tau\Big(\sup_\alpha A_\alpha\Big) = \sup_\alpha \tau(A_\alpha),
	\]
	for every bounded increasing net $(A_\alpha)$ of elements of $M_+$,
	with the supremum on the left taken in the strong operator topology of
	Definition~\ref{def:wot-sot}.
\end{definition}

\begin{definition}[Semifinite trace]\label{def:semifinite-trace}
	Let $\tau$ be a trace on $M$, as per Definition~\ref{def:trace-functional}.
	The trace $\tau$ is semifinite if, for every nonzero $A \in M_+$, there
	exists a nonzero $B \in M_+$ with $B \leq A$, as per
	Remark~\ref{rem:projector-order} extended to the order induced on all
	of $M_+$, and $\tau(B) < \infty$.
\end{definition}

\begin{remark}
	Definition~\ref{def:normal-trace} is the same continuity condition as
	Definition~\ref{def:normal-state}, extended from normalized states to
	the possibly unbounded functional $\tau$. When $M = B(H)$ and $\tau =
	\mathrm{Tr}$, as per Definition~\ref{def:trace}, restricted to the
	whole of $B(H)_+$ rather than to $B_1(H)_+$, $\tau$ is faithful, normal,
	and semifinite, and every density matrix $\rho$ of Definition 2.3.15
	satisfies $\mathrm{Tr}(\rho) < \infty$, consistently with
	Definition~\ref{def:semifinite-trace} applied to $A = 1$.
\end{remark}

With finiteness of projections, factoriality, and the trace all in
place, the classification of Murray and von Neumann can now be stated.
It shows that the existence of a trace on a factor is entirely
determined by the combinatorics of its projections, so that the
absence of minimal projections already forces a choice between two
mutually exclusive possibilities for the whole algebra.

\begin{theorem}[Murray-von Neumann classification of factors]\label{thm:murray-von-neumann}
	Let $M$ be a factor on $H$, as per Definition~\ref{def:factor}. Exactly
	one of the following holds.
	\begin{enumerate}
		\item Type I: $M$ has a minimal nonzero projector, that is, a
		projector $P \in M$, as per Definition 2.2.17, such that no projector
		$Q \in M$ satisfies $0 \neq Q < P$ as per
		Remark~\ref{rem:projector-order}. In this case $M$ is $*$-isomorphic,
		as per Definition 2.3.1, to $B(K)$ for some Hilbert space $K$.
		\item Type II: $M$ has no minimal nonzero projector, but $M$ admits a
		faithful normal semifinite trace, as per Definitions
		\ref{def:trace-functional}, \ref{def:faithful-trace},
		\ref{def:normal-trace}, and \ref{def:semifinite-trace}.
		\item Type III: $M$ has no minimal nonzero projector, and $M$ admits
		no nonzero faithful normal semifinite trace.
	\end{enumerate}
\end{theorem}
\noindent This is proved in \cite[TODO: collocazione precisa]{TODO-MurrayVonNeumann}.

\begin{remark}
	Type III factors are not a single, further indistinguishable class.
	This third case of Theorem~\ref{thm:murray-von-neumann} splits into a
	family of mutually exclusive subtypes, denoted $\mathrm{III}_0$,
	$\mathrm{III}_\lambda$ for $\lambda \in (0,1)$, and $\mathrm{III}_1$,
	detected by an invariant built from the modular operator only once the
	tools of Section 4.4 are available, and Section 4.9 returns to this
	finer classification in full. \cite[TODO: collocazione
	precisa]{TODO-fineTypeIII}
\end{remark}

The type III case of Theorem~\ref{thm:murray-von-neumann} is where the
framework of Chapter 2 breaks down completely, and it is worth spelling
out exactly what is lost.

\begin{remark}[Collapse of the density matrix picture]\label{rem:type-III-density-matrix}
	Let $M$ be a factor of type III on $H$, as per
	Theorem~\ref{thm:murray-von-neumann}. Since $M$ admits no nonzero
	faithful normal semifinite trace, no normal state $\omega$ on $M$, as
	per Definition~\ref{def:normal-state}, can be represented as
	$\omega(A) = \tau(\rho A)$ for a fixed positive $\rho \in M$ with
	respect to any trace $\tau$ on $M$, in contrast with the description of
	Definition 2.3.15 and Proposition 2.3.14 available for $B(H)$ itself.
	Every normal state on $M$ is nevertheless still a vector state for some
	faithful representation of $M$, by Remark~\ref{rem:vector-normal}
	together with the Gelfand-Naimark-Segal construction of Theorem 2.7,
	so a Hilbert space description of the state survives even though a
	density matrix does not. Algebras of local observables in relativistic
	quantum field theory are type III on physical grounds, so that this
	collapse is the generic situation for such algebras rather than a
	pathological exception. \cite[TODO: collocazione precisa]{TODO-QFTtypeIII}
\end{remark}

Theorem~\ref{thm:murray-von-neumann} therefore leaves this thesis with
a genuine gap. Chapter 2 built the entire probabilistic content of
quantum mechanics, states as density matrices and time evolution as
conjugation by $e^{-itH/\hbar}$, on the trace of $B(H)$, a type I
factor by Theorem~\ref{thm:murray-von-neumann} applied with $M = B(H)$
itself, since every rank-one projector is minimal there. On a type III
factor, neither notion is available in that form. The remainder of
this chapter constructs, for an arbitrary von Neumann algebra $M$
together with a vector cyclic and separating for it as in
Lemma~\ref{lem:cyclic-separating}, a canonical one-parameter group of
automorphisms of $M$ that requires no trace whatsoever, built instead
directly from the vector state itself through Tomita's theorem of
Section 4.4. Section 4.6 will verify explicitly that, when $M$ does
possess a trace, this group reduces to ordinary Gibbs time evolution,
so that the construction to follow genuinely extends, rather than
replaces, the familiar type I picture of Chapter 2.

\section{The Tomita--Takesaki Theorem for $M$}\label{sec:TT-for-M}

Section 4.2 exhibited von Neumann algebras on which no trace exists, so the symmetrized construction underlying Definition 1.6.23 and the Gibbs dynamics it generalizes cannot be the route to a canonical time-like flow on every von Neumann algebra. Section 4.3 developed, with no algebra present, a replacement built entirely from the geometry of a single closed real subspace of a Hilbert space: a standard subspace $K$ determines, through its Tomita operator and the resulting polar decomposition, a modular operator $\Delta_K$ and a modular conjugation $J_K$, and Theorem~\ref{thm:fundamental-standard} showed that $\Delta_K^{it}$ leaves $K$ invariant for every real $t$, while $J_K$ exchanges $K$ with its symplectic complement $K'$. This section supplies the missing link between that abstract theory and the von Neumann algebra $M$ of the preceding sections. Given $M$ together with a cyclic and separating vector $\Omega$, as per Lemma~\ref{lem:cyclic-separating}, the real closed subspace generated by the symmetrized vectors $A\Omega + A^*\Omega$ for $A$ ranging over $M$ turns out to be standard in the sense of Definition~\ref{def:standard-subspace}, and its symplectic complement coincides with the analogous subspace built from the commutant $M'$. Feeding this identification into Theorem~\ref{thm:fundamental-standard} promotes the modular operator and modular conjugation of this subspace into a one-parameter group of automorphisms of $M$ and an antiunitary exchanging $M$ with $M'$, with neither object referring to a trace at any point of the construction. This is the Tomita--Takesaki theorem, and the one-parameter group it produces is the object that Section 4.5 onward treats as the canonical substitute for time evolution on an arbitrary von Neumann algebra.

\begin{definition}[The real subspace generated by $M$]\label{def:KM-subspace}
	Let $M$ be a von Neumann algebra on $H$ with cyclic and separating vector $\Omega$, as per Lemma~\ref{lem:cyclic-separating}. The real subspace generated by $M$ is
	\[
	K_M := \overline{\operatorname{span}_{\mathbb{R}} \{A\Omega + A^*\Omega \mid A \in M\}},
	\]
	where the closure and the real span are taken in the norm topology of $H$.
\end{definition}

Before $K_M$ can be fed into the apparatus of Section 4.3, it has to be checked against the two conditions of Definition~\ref{def:standard-subspace}. The first of these, density of $K_M + iK_M$, follows from cyclicity of $\Omega$ alone and an elementary identity relating $A\Omega$, for $A \in M$, to the two generators $A\Omega + A^*\Omega$ and $(iA)\Omega + (iA)^*\Omega$.

\begin{lemma}[$K_M + iK_M$ is dense]\label{lem:KM-density}
	Let $K_M$ be as per Definition~\ref{def:KM-subspace}. Then $\overline{K_M + iK_M} = H$.
\end{lemma}

\begin{proof}
	Fix $A \in M$ and set $\xi_A := A\Omega + A^*\Omega$ and $\xi_{iA} := (iA)\Omega + (iA)^*\Omega$, both lying in $K_M$ since $iA \in M$ as well, $M$ being a complex algebra. Expanding the second generator using antilinearity of the involution yields
	\[
	\xi_{iA} = iA\Omega + \overline{i}\,A^*\Omega = iA\Omega - iA^*\Omega = i(A\Omega - A^*\Omega).
	\]
	Solving the resulting linear system
	\[
	A\Omega + A^*\Omega = \xi_A, \qquad A\Omega - A^*\Omega = -i\xi_{iA},
	\]
	for $A\Omega$ gives
	\[
	A\Omega = \tfrac{1}{2}\xi_A + i\left(-\tfrac{1}{2}\xi_{iA}\right).
	\]
	Since $K_M$ is real-linear and contains $\xi_A$ and $\xi_{iA}$, both $\tfrac{1}{2}\xi_A$ and $-\tfrac{1}{2}\xi_{iA}$ lie in $K_M$, so $A\Omega \in K_M + iK_M$. As $A \in M$ was arbitrary, $M\Omega \subseteq K_M + iK_M$, and $M\Omega$ is dense in $H$ because $\Omega$ is cyclic for $M$. It descends that $K_M + iK_M$ is dense as well.
\end{proof}

The second condition of Definition~\ref{def:standard-subspace}, that $K_M \cap iK_M$ reduces to the zero vector, is considerably less direct, and is obtained here by locating $K_M$ inside the symplectic complement of the analogous subspace $K_{M'}$ built from the commutant, together with a general fact about symplectic complements that holds for an arbitrary closed real subspace of a Hilbert space, with no algebra involved.

\begin{lemma}[Symplectic self-intersection and orthogonality]\label{lem:complement-intersection}
	Let $K \subseteq H$ be a closed real-linear subspace, with symplectic complement $K'$, as per Definition~\ref{def:symplectic-complement}. Then
	\[
	K' \cap iK' = \overline{K+iK}^{\perp},
	\]
	the orthogonal complement, with respect to the inner product of $H$, of the closure of $K+iK$.
\end{lemma}

\begin{proof}
	Fix $\xi \in H$. By Definition~\ref{def:symplectic-complement}, $\xi \in K'$ if and only if $\mathrm{Im}(\xi,\eta) = 0$ for every $\eta \in K$. Antilinearity of the inner product in its first argument, as per Definition 2.2.1, gives $(-i\xi,\eta) = \overline{-i}\,(\xi,\eta) = i(\xi,\eta)$, so that $\mathrm{Im}(-i\xi,\eta) = \mathrm{Re}(\xi,\eta)$. Hence $\xi \in iK'$, that is, $-i\xi \in K'$, if and only if $\mathrm{Re}(\xi,\eta) = 0$ for every $\eta \in K$. Combining the two characterizations, $\xi \in K' \cap iK'$ if and only if $(\xi,\eta) = 0$ for every $\eta \in K$. Linearity of the inner product in its second argument then gives, for $\eta,\eta' \in K$,
	\[
	(\xi,\eta+i\eta') = (\xi,\eta) + i(\xi,\eta') = 0,
	\]
	so that $\xi$ orthogonal to every element of $K$ is equivalent to $\xi$ orthogonal to every element of $K+iK$, and hence, by continuity of the inner product, to every element of $\overline{K+iK}$. It descends that $K'\cap iK' = \overline{K+iK}^{\perp}$.
\end{proof}

Lemma~\ref{lem:complement-intersection} makes no reference to standardness of $K$, and is applied below to $K_{M'}$ rather than to $K_M$ itself. What remains is to locate $K_{M'}$ inside $K_M'$, equivalently $K_M$ inside $K_{M'}'$, the two inclusions being the same statement read in opposite order. This is a direct computation with the inner product of $H$, using only that $M$ and $M'$ commute elementwise.

\begin{lemma}[Symplectic orthogonality of $M$ and $M'$]\label{lem:MMp-orthogonal}
	For every $A \in M$ and $A' \in M'$,
	\[
	\mathrm{Im}\big(A\Omega + A^*\Omega,\ A'\Omega + A'^*\Omega\big) = 0.
	\]
\end{lemma}

\begin{proof}
	Write $x := A\Omega + A^*\Omega$ and $y := A'\Omega + A'^*\Omega$. Using Proposition 2.2.21 to move each operator from the first to the second argument of the inner product, $(B\phi,\psi) = (\phi,B^*\psi)$ for every bounded $B$, each of the four terms in the expansion of $(x,y)$ reduces to a form with no operator acting on the first argument:
	\[
	(x,y) = (\Omega,A^*A'\Omega) + (\Omega,A^*A'^*\Omega) + (\Omega,AA'\Omega) + (\Omega,AA'^*\Omega) = \big(\Omega,(A+A^*)(A'+A'^*)\Omega\big).
	\]
	The identical computation with the roles of $A$ and $A'$ exchanged gives $(y,x) = \big(\Omega,(A'+A'^*)(A+A^*)\Omega\big)$. Since $A+A^* \in M$ and $A'+A'^* \in M'$, and $M$ and $M'$ commute elementwise by definition of the commutant, $(A+A^*)(A'+A'^*) = (A'+A'^*)(A+A^*)$, so $(x,y) = (y,x)$. Conjugate symmetry of the inner product, as per Definition 2.2.1, gives $(y,x) = \overline{(x,y)}$, so $(x,y) = \overline{(x,y)}$, meaning $(x,y) \in \mathbb{R}$ and $\mathrm{Im}(x,y) = 0$.
\end{proof}

\noindent Extending Lemma~\ref{lem:MMp-orthogonal} from the two generating sets to the subspaces they span and close is immediate, since $\mathrm{Im}(\cdot,\cdot)$ is a continuous real-bilinear form and both $K_M$ and $K_{M'}$ are real-linear and closed.

\begin{corollary}\label{cor:KM-in-KMp-complement}
	$K_M \subseteq K_{M'}'$ and $K_{M'} \subseteq K_M'$.
\end{corollary}

\begin{proof}
	By Lemma~\ref{lem:MMp-orthogonal}, every generator of $K_M$ has vanishing symplectic product with every generator of $K_{M'}$, as per Definition~\ref{def:symplectic-complement}. Real-bilinearity and continuity of $\mathrm{Im}(\cdot,\cdot)$ extend this to all of $K_M$ against all of $K_{M'}$, giving $K_M \subseteq K_{M'}'$. The second inclusion is the same statement with the two subspaces exchanged.
\end{proof}

Corollary~\ref{cor:KM-in-KMp-complement} and Lemma~\ref{lem:complement-intersection} combine directly into the standardness of $K_M$, once the companion subspace $K_{M'}$ is checked against the density half of Definition~\ref{def:standard-subspace}, which holds by the same argument as Lemma~\ref{lem:KM-density} applied to $M'$ in place of $M$, since $\Omega$ is cyclic for $M'$ as well as for $M$ by Lemma~\ref{lem:cyclic-separating}.

\begin{proposition}[$K_M$ is standard]\label{prop:KM-standard}
	$K_M$, as per Definition~\ref{def:KM-subspace}, is a standard subspace of $H$, as per Definition~\ref{def:standard-subspace}.
\end{proposition}

\begin{proof}
	Density of $K_M + iK_M$ is Lemma~\ref{lem:KM-density}. For the intersection condition, Corollary~\ref{cor:KM-in-KMp-complement} gives $K_M \subseteq K_{M'}'$, hence $K_M \cap iK_M \subseteq K_{M'}' \cap iK_{M'}'$. By Lemma~\ref{lem:complement-intersection} applied to $K = K_{M'}$,
	\[
	K_{M'}' \cap iK_{M'}' = \overline{K_{M'} + iK_{M'}}^{\perp}.
	\]
	Since $\Omega$ is cyclic for $M'$ as well as for $M$, by Lemma~\ref{lem:cyclic-separating}, the argument of Lemma~\ref{lem:KM-density} applies verbatim with $M'$ in place of $M$, giving $\overline{K_{M'}+iK_{M'}} = H$ and therefore $\overline{K_{M'}+iK_{M'}}^{\perp} = \{0\}$. It descends that $K_M \cap iK_M = \{0\}$, so $K_M$ satisfies both conditions of Definition~\ref{def:standard-subspace}.
\end{proof}

Proposition~\ref{prop:KM-standard} unlocks the full strength of Theorem~\ref{thm:fundamental-standard} for $K_M$, in particular the identity $K_M'' = K_M$ quoted in the remark following Definition~\ref{def:symplectic-complement}. Corollary~\ref{cor:KM-in-KMp-complement} already supplies half of the identification between $K_M'$ and $K_{M'}$. The reverse inclusion is the actual commutation theorem of Tomita--Takesaki theory, the one genuinely deep fact about $K_M$ and $K_{M'}$ not already contained in the elementary computation of Lemma~\ref{lem:MMp-orthogonal}, and is quoted here from the literature rather than reproduced.

\begin{proposition}[$K_M' = K_{M'}$]\label{prop:KM-complement}
	$K_M' = K_{M'}$.
\end{proposition}

\begin{proof}
	Corollary~\ref{cor:KM-in-KMp-complement} gives $K_{M'} \subseteq K_M'$. The reverse inclusion $K_M' \subseteq K_{M'}$ is proved in \cite[TODO: collocazione precisa]{TODO-RieffelVanDaele}, in the equivalent formulation of \cite[TODO: collocazione precisa]{TODO-Longo}.
\end{proof}

With $K_M$ standard and its symplectic complement identified, Section 4.3 attaches to it a modular operator and a modular conjugation exactly as it does to any standard subspace. These two operators are what the remainder of this section turns into the modular automorphism group of $M$ itself.

\begin{definition}[Modular operator and modular conjugation of $M$]\label{def:modular-M}
	Let $M$ be a von Neumann algebra on $H$ with cyclic and separating vector $\Omega$, and let $K_M$ be as per Definition~\ref{def:KM-subspace}, standard by Proposition~\ref{prop:KM-standard}. The modular operator of $M$ relative to $\Omega$ is $\Delta := \Delta_{K_M}$, and the modular conjugation of $M$ relative to $\Omega$ is $J := J_{K_M}$, both as per Definition~\ref{def:modular-K}.
\end{definition}

Before $\Delta$ and $J$ are shown to implement the Tomita--Takesaki theorem on $M$, it is worth checking that the Tomita operator $S_{K_M}$ of Definition~\ref{def:tomita-K}, defined on all of $K_M + iK_M$ through the real subspace $K_M$, reduces on the smaller domain $M\Omega$ to the antilinear map $A\Omega \mapsto A^*\Omega$ usually taken as the starting point of modular theory. This identification is what justifies calling $\Delta$ and $J$ the modular operator and modular conjugation of $M$, rather than merely of the auxiliary subspace $K_M$.

\begin{lemma}\label{lem:SKM-on-MOmega}
	For every $A \in M$, $A\Omega \in K_M + iK_M$ and $S_{K_M}(A\Omega) = A^*\Omega$.
\end{lemma}

\begin{proof}
	Membership $A\Omega \in K_M + iK_M$, together with the decomposition $A\Omega = \tfrac{1}{2}\xi_A + i(-\tfrac{1}{2}\xi_{iA})$, is Lemma~\ref{lem:KM-density}. Applying Definition~\ref{def:tomita-K} to this decomposition,
	\[
	S_{K_M}(A\Omega) = \tfrac{1}{2}\xi_A - i\left(-\tfrac{1}{2}\xi_{iA}\right) = \tfrac{1}{2}(A\Omega+A^*\Omega) + \tfrac{i}{2}\xi_{iA}.
	\]
	Substituting $\xi_{iA} = i(A\Omega - A^*\Omega)$, as per the proof of Lemma~\ref{lem:KM-density}, into the second term gives $\tfrac{i}{2}\xi_{iA} = \tfrac{i^2}{2}(A\Omega-A^*\Omega) = -\tfrac{1}{2}(A\Omega-A^*\Omega)$, so
	\[
	S_{K_M}(A\Omega) = \tfrac{1}{2}(A\Omega+A^*\Omega) - \tfrac{1}{2}(A\Omega-A^*\Omega) = A^*\Omega. \qedhere
	\]
\end{proof}

Lemma~\ref{lem:SKM-on-MOmega} identifies $\Delta$ and $J$ as the natural extensions of the algebraic operations $A\Omega \mapsto A\Omega$ and $A\Omega \mapsto A^*\Omega$ on $M$, so that Theorem~\ref{thm:fundamental-standard}, already available for $K_M$ by Propositions~\ref{prop:KM-standard} and~\ref{prop:KM-complement}, can be read as a statement about $M$ once the invariance it grants the subspace $K_M$ is promoted to an invariance of the algebra $M$.

\begin{theorem}[Tomita--Takesaki theorem]\label{thm:TT-M}
	Let $M$ be a von Neumann algebra on $H$ with cyclic and separating vector $\Omega$, as per Lemma~\ref{lem:cyclic-separating}, and let $\Delta$, $J$ be its modular operator and modular conjugation relative to $\Omega$, as per Definition~\ref{def:modular-M}. Then
	\[
	\Delta^{it} M \Delta^{-it} = M \quad \forall t \in \mathbb{R}, \qquad J M J = M'.
	\]
\end{theorem}

\begin{proof}
	Propositions~\ref{prop:KM-standard} and~\ref{prop:KM-complement} place $K_M$ within the hypotheses of Theorem~\ref{thm:fundamental-standard}, so part 3 of that theorem applies directly to $K = K_M$ and gives
	\[
	\Delta^{it} K_M = K_M \quad \forall t \in \mathbb{R}, \qquad J K_M = K_M' = K_{M'},
	\]
	the second equality by Proposition~\ref{prop:KM-complement}. This fixes the action of $\Delta^{it}$ and $J$ on the real subspace generated by the symmetrized vectors $A\Omega + A^*\Omega$, but the stated theorem asserts an invariance of the algebra $M$ itself under conjugation by $\Delta^{it}$ and $J$, a strictly stronger statement controlling $\Delta^{it}A\Delta^{-it}$ and $JAJ$ for each individual $A \in M$, rather than only the action on symmetrized vectors. Promoting the subspace-level invariance just established to this operator-level statement is proved in \cite[TODO: collocazione precisa]{TODO-RieffelVanDaele}.
\end{proof}

\noindent The one-parameter group of automorphisms this theorem attaches to $M$ depends on the choice of $\Omega$ only through the vector itself, with no auxiliary data, and is named accordingly.

\begin{definition}[Modular automorphism group]\label{def:modular-automorphism}
	Let $M$, $\Omega$, $\Delta$ be as per Theorem~\ref{thm:TT-M}. The modular automorphism group of $M$ relative to $\Omega$ is the one-parameter group $\sigma^\Omega = (\sigma^\Omega_t)_{t \in \mathbb{R}}$ of $*$-automorphisms of $M$ defined by
	\[
	\sigma^\Omega_t(A) := \Delta^{it} A \Delta^{-it}, \qquad A \in M.
	\]
\end{definition}

\begin{remark}
	Definition~\ref{def:modular-automorphism} produces, for every von Neumann algebra $M$ with a cyclic and separating vector, a canonical one-parameter group of $*$-automorphisms of $M$, with no appeal to a trace, a Hamiltonian, or any externally specified dynamics, in contrast with the time evolution of Definition 1.6.23, which is simply postulated as additional data on top of the commutative algebra $A$. The notation $\sigma^\Omega$ flags that the group is defined directly from the vector $\Omega$, not yet from the state it induces; Section 4.7 shows that $\sigma^\Omega_t$ in fact depends on $\Omega$ only through the state $\omega(\cdot) := (\Omega,\cdot\,\Omega)$, which justifies writing $\sigma^\omega_t$ from that point onward. Whether $\sigma^\Omega$ deserves to be read as a dynamics in the physical sense asked of Definition 1.6.23, rather than as a piece of purely geometric bookkeeping attached to $\Omega$, is the question taken up in Section 4.5 through the Kubo--Martin--Schwinger condition.
\end{remark}

\section{Example IV: Modular Theory of a Gibbs State on $M_n(\mathbb{C})$}\label{sec:example-gibbs}

This example instantiates the Tomita--Takesaki theorem of Section 4.4 and the Kubo-Martin-Schwinger characterization of Section 4.5 on the simplest possible algebra, the full matrix algebra $M_n(\mathbb{C})$, already identified in Section 4.2 as a type I factor carrying a trace. Fixing a Gibbs state on $M_n(\mathbb{C})$ makes every object constructed abstractly in the preceding two sections, the modular operator, the modular conjugation, and the modular automorphism group, computable in closed form, with no unbounded operator subtlety and no appeal to the literature anywhere in the computation. Working through this example also closes a thread left open in Section 2.7: the maximally mixed state $\omega_0$ of the quantum spin system, there purified by hand into the vector $\Psi_{\omega_0}$, reappears here as the $\beta \to 0$ limit of the Gibbs state, with the modular operator collapsing to the identity exactly as a trace-induced construction should.

\paragraph{The Gibbs state and its GNS representation.}
Fix $n \in \mathbb{N}$ and a selfadjoint $H \in M_n(\mathbb{C})$, the Hamiltonian, and for $\beta > 0$ set $Z_\beta := \mathrm{Tr}(e^{-\beta H})$, strictly positive since $e^{-\beta H}$ is positive and invertible. The Gibbs state at inverse temperature $\beta$ is
\[
\omega_\beta(A) := \frac{1}{Z_\beta}\mathrm{Tr}\big(e^{-\beta H}A\big), \qquad A \in M_n(\mathbb{C}).
\]
$M_n(\mathbb{C})$ is already a von Neumann algebra on $\mathbb{C}^n$ in its defining representation, with $M_n(\mathbb{C})' = \mathbb{C}1$ and trivial center, the elementary instance of the type I factor of Section 4.2. The construction of this example represents $M_n(\mathbb{C})$ not on $\mathbb{C}^n$ but on the Hilbert space of $M_n(\mathbb{C})$ itself, equipped with the Hilbert-Schmidt inner product $(X,Y) := \mathrm{Tr}(X^*Y)$, as per Example 2.2.4, so write $H_{\omega_\beta} := M_n(\mathbb{C})$ for this Hilbert space to keep it notationally distinct from the algebra acting on it. Let $\pi_{\omega_\beta} : M_n(\mathbb{C}) \to B(H_{\omega_\beta})$ be left multiplication, $\pi_{\omega_\beta}(A)X := AX$, and set $M := \pi_{\omega_\beta}(M_n(\mathbb{C}))$. Since $H_{\omega_\beta}$ is finite-dimensional, every subalgebra of $B(H_{\omega_\beta})$ is automatically weakly closed, so $M$ is a von Neumann algebra on $H_{\omega_\beta}$, $*$-isomorphic to $M_n(\mathbb{C})$ and hence itself a type I factor carrying the trace $\mathrm{Tr}$.

A direct computation identifies the commutant of $M$. If $T \in B(H_{\omega_\beta})$ commutes with $\pi_{\omega_\beta}(A)$ for every $A \in M_n(\mathbb{C})$, setting $B := T(1)$ gives, for every $A \in M_n(\mathbb{C})$,
\[
T(A) = T(A \cdot 1) = T\big(\pi_{\omega_\beta}(A)1\big) = \pi_{\omega_\beta}(A)T(1) = AB,
\]
so $T$ is right multiplication by $B$. Conversely every right multiplication commutes with every left multiplication. Hence $M' = \{R_B : B \in M_n(\mathbb{C})\}$, where $R_B(X) := XB$.

The cyclic and separating vector representing $\omega_\beta$ is the normalized square root of the density matrix,
\[
\Omega_\beta := \frac{1}{\sqrt{Z_\beta}}e^{-\beta H/2} \in H_{\omega_\beta}.
\]
It has unit norm, since $(\Omega_\beta,\Omega_\beta) = Z_\beta^{-1}\mathrm{Tr}(e^{-\beta H/2}e^{-\beta H/2}) = Z_\beta^{-1}\mathrm{Tr}(e^{-\beta H}) = 1$, using that $e^{-\beta H/2}$ is selfadjoint. It induces $\omega_\beta$ as a vector state, since for $A \in M_n(\mathbb{C})$,
\[
(\Omega_\beta,\pi_{\omega_\beta}(A)\Omega_\beta) = \frac{1}{Z_\beta}\mathrm{Tr}\big(e^{-\beta H/2}Ae^{-\beta H/2}\big) = \frac{1}{Z_\beta}\mathrm{Tr}\big(e^{-\beta H}A\big) = \omega_\beta(A),
\]
the middle equality by cyclicity of the trace. Since $e^{-\beta H/2}$ is invertible, the map $A \mapsto Ae^{-\beta H/2}$ is a bijection of $M_n(\mathbb{C})$ onto itself, so $\pi_{\omega_\beta}(M_n(\mathbb{C}))\Omega_\beta = H_{\omega_\beta}$, and $\Omega_\beta$ is cyclic for $M$. For the same reason $\pi_{\omega_\beta}(A)\Omega_\beta = 0$ forces $Ae^{-\beta H/2}=0$ and hence $A=0$, so $\Omega_\beta$ is separating for $M$ as well.

\paragraph{The modular operator and modular conjugation.}
By Lemma~\ref{lem:SKM-on-MOmega}, the Tomita operator of $M$ relative to $\Omega_\beta$ satisfies $S(\pi_{\omega_\beta}(A)\Omega_\beta) = \pi_{\omega_\beta}(A^*)\Omega_\beta$ for every $A \in M_n(\mathbb{C})$. Writing a general $X \in H_{\omega_\beta}$ as $X = \pi_{\omega_\beta}(A)\Omega_\beta$ with $A = \sqrt{Z_\beta}\,Xe^{\beta H/2}$, the unique preimage since $e^{-\beta H/2}$ is invertible, gives
\begin{equation}\label{eq:S-gibbs}
	S(X) = e^{\beta H/2}X^*e^{-\beta H/2}, \qquad X \in H_{\omega_\beta}.
\end{equation}
Set $J(X) := X^*$ and $\Delta^{1/2}(X) := e^{-\beta H/2}Xe^{\beta H/2}$. $J$ is an antiunitary involution, since $\|X^*\|^2 = \mathrm{Tr}(XX^*) = \mathrm{Tr}(X^*X) = \|X\|^2$ and $(X^*)^*=X$, and $\Delta^{1/2}$ is positive and selfadjoint, being the product of the commuting positive operators $L_{e^{-\beta H/2}}$ and $R_{e^{\beta H/2}}$, each positive because $\mathrm{Tr}\big(X^*e^{-\beta H/2}X\big) = \|e^{-\beta H/4}Xe^{-\beta H/4}\|^2 \geq 0$ and similarly for $R_{e^{\beta H/2}}$. A direct check confirms the polar decomposition,
\[
J\Delta^{1/2}(X) = J\big(e^{-\beta H/2}Xe^{\beta H/2}\big) = \big(e^{-\beta H/2}Xe^{\beta H/2}\big)^* = e^{\beta H/2}X^*e^{-\beta H/2} = S(X),
\]
using that $e^{\pm\beta H/2}$ is selfadjoint. Since $S$ is invertible on the finite-dimensional space $H_{\omega_\beta}$, being an involution, its polar decomposition is unique, so $J$ and $\Delta^{1/2}$ as just defined are the modular conjugation and the square root of the modular operator of $M$ relative to $\Omega_\beta$, as per Definition~\ref{def:modular-M}. Squaring $\Delta^{1/2}$ and writing $\rho := e^{-\beta H}$ for the density matrix,
\begin{equation}\label{eq:Delta-gibbs}
	\Delta(X) = \rho X \rho^{-1}, \qquad J(X) = X^*, \qquad X \in H_{\omega_\beta}.
\end{equation}

\paragraph{The modular automorphism group.}
Since $L_\rho$ and $R_{\rho^{-1}}$ commute, $\Delta^{it} = L_{\rho^{it}}R_{\rho^{-it}}$, that is, $\Delta^{it}(X) = \rho^{it}X\rho^{-it}$. For $A \in M_n(\mathbb{C})$ and $X \in H_{\omega_\beta}$,
\[
\Delta^{it}\pi_{\omega_\beta}(A)\Delta^{-it}(X) = \Delta^{it}\big(A\rho^{-it}X\rho^{it}\big) = \rho^{it}A\rho^{-it}\cdot X \cdot \rho^{it}\rho^{-it} = \big(\rho^{it}A\rho^{-it}\big)X,
\]
so the modular automorphism group of Definition~\ref{def:modular-automorphism} acts on $M_n(\mathbb{C})$ itself as
\begin{equation}\label{eq:sigma-gibbs}
	\sigma_t^{\Omega_\beta}(A) = \rho^{it}A\rho^{-it} = e^{-i\beta t H}Ae^{i\beta t H}, \qquad A \in M_n(\mathbb{C}).
\end{equation}
Since $H$ is selfadjoint, $-i\beta t H$ is skew-adjoint for every $t \in \mathbb{R}$, so $\rho^{it}=e^{-i\beta t H}$ is unitary and \eqref{eq:sigma-gibbs} is manifestly an automorphism of $M_n(\mathbb{C})$, confirming $\Delta^{it}M\Delta^{-it}=M$ of Theorem~\ref{thm:TT-M} directly, with no appeal to the cited step of its proof. Likewise $J\pi_{\omega_\beta}(A)J(X) = J(AX^*) = XA^* = R_{A^*}(X)$, so $J\pi_{\omega_\beta}(A)J = R_{A^*} \in M'$, and as $A$ ranges over $M_n(\mathbb{C})$ so does $A^*$, confirming $JMJ=M'$ concretely as well.

Equation~\eqref{eq:sigma-gibbs} exhibits $\sigma^{\Omega_\beta}_t$ as ordinary Heisenberg evolution under $H$ at physical time $\tau = -\beta t$, rather than at modular time $t$ itself. This rescaling is why Theorem~\ref{thm:takesaki} places $\omega_\beta$ at Kubo-Martin-Schwinger value exactly $1$ with respect to $\sigma^{\Omega_\beta}$: the inverse temperature $\beta$ of the Gibbs state is already absorbed into the definition of $\Omega_\beta$, and with it into $\Delta$, so that modular time is dimensionless and reproduces equilibrium at the physical inverse temperature $\beta$ only once rescaled by $\beta$ back into ordinary time. Section 4.7 returns to this rescaling as the central feature of the thermal time hypothesis.

\begin{remark}[The limit $\beta \to 0$]
	As $\beta \to 0$, $\rho = e^{-\beta H} \to 1$, so \eqref{eq:Delta-gibbs} and \eqref{eq:sigma-gibbs} give $\Delta \to 1$ and $\sigma_t^{\Omega_0} = \mathrm{id}$ for every $t$, while $J$ remains the matrix adjoint $X \mapsto X^*$ regardless of $\beta$. The state $\omega_0(A) = \mathrm{Tr}(A)/n$ is the maximally mixed state, tracial on $M_n(\mathbb{C})$, and $\Omega_0 = 1/\sqrt{n}$, the identity matrix of $H_{\omega_\beta}$ rescaled. Under the matrix-unit isomorphism $M_n(\mathbb{C}) \cong \mathbb{C}^n \otimes \mathbb{C}^n$, $E_{ij} \leftrightarrow e_i \otimes e_j$, used in Section 2.7, the identity matrix corresponds to $\sum_k e_k \otimes e_k$, so $\Omega_0 \leftrightarrow n^{-1/2}\sum_k e_k \otimes e_k$, which for $n=2$ is exactly the vector $\Psi_{\omega_0}$ of Equation (2.67). The collapse of $\Delta$ to the identity precisely when $\omega$ is tracial is not a coincidence of this example, and is returned to in general in Section 4.10.
\end{remark}

\section{The Thermal Time Hypothesis and the Cocycle Theorem of Connes}\label{sec:thermal-time}

Theorem~\ref{thm:takesaki} singled out $\sigma^\Omega$, among all one-parameter groups of automorphisms of $M$, as the unique one compatible with equilibrium for $\omega$. Connes and Rovelli proposed to invert this reading: rather than starting from a Hamiltonian and verifying that the resulting Gibbs state is $\sigma^\Omega$-KMS, start from a physical system with no privileged notion of time at all, described purely by a von Neumann algebra of observables $M$ and whatever state $\omega$ the system happens to be found in, and declare the modular automorphism group $\sigma^\omega$ itself to be the system's time evolution. For this proposal to deserve the name of a physical hypothesis, rather than an arbitrary choice tied to the particular state $\omega$, the flow it produces must not depend essentially on which state was used to define it. This section makes that requirement precise and proves it, via the two-by-two matrix construction of Connes, which embeds any two states of $M$ into a single state of a larger algebra and reads off, from the modular theory of the larger algebra, a unitary cocycle relating the two original modular flows.

\begin{definition}[Thermal time hypothesis]\label{def:thermal-time}
	Let $M$ be a von Neumann algebra describing a physical system, with no observable of $M$ singled out as a Hamiltonian. The thermal time hypothesis asserts that the physical time evolution of the system, once it is found in a faithful normal state $\omega$, is the modular automorphism group $\sigma^\omega$ of Definition~\ref{def:modular-automorphism}, and that this evolution is a genuine physical fact about the system rather than an artifact of the choice of $\omega$ among the states the system might occupy.
\end{definition}

Taken literally, Definition~\ref{def:thermal-time} requires $\sigma^\omega$ and $\sigma^{\omega'}$, for two different faithful normal states $\omega,\omega'$ on the same $M$, to coincide as physical evolutions whenever they coincide as the automorphisms that matter physically, that is, up to conjugation by a unitary already present in $M$. Every such conjugation is an inner automorphism, and inner automorphisms of the observable algebra correspond to no observable change of description.

\begin{definition}[Inner and outer automorphisms]\label{def:inn-out}
	Let $M$ be a von Neumann algebra. For a unitary $u \in M$, write $\mathrm{Ad}(u)$ for the automorphism $\mathrm{Ad}(u)(A) := uAu^*$, $A \in M$. The group $\mathrm{Inn}(M) := \{\mathrm{Ad}(u) \mid u \in M \text{ unitary}\}$ is the group of inner automorphisms of $M$, and $\mathrm{Out}(M) := \mathrm{Aut}(M)/\mathrm{Inn}(M)$ is the group of outer automorphisms.
\end{definition}

The content of this section is the proof that $\sigma_t^\varphi$ and $\sigma_t^\psi$ represent the same element of $\mathrm{Out}(M)$ for every $t$ and every pair of faithful normal states $\varphi,\psi$ on $M$, Theorem~\ref{thm:out-independence} below, which makes Definition~\ref{def:thermal-time} well posed. The route to it passes through an explicit unitary, the Connes cocycle, relating $\sigma^\varphi$ and $\sigma^\psi$ directly.

\paragraph{The balanced algebra and the balanced state.}
Fix a von Neumann algebra $M$ on a Hilbert space $H$ and two faithful normal states $\varphi,\psi$ on $M$. Let $N$ be the set of $2\times2$ matrices with entries in $M$,
\[
N := \left\{ X = \begin{pmatrix} A_{11} & A_{12} \\ A_{21} & A_{22} \end{pmatrix} \ \middle|\ A_{ij} \in M \right\},
\]
with entrywise matrix multiplication and involution $X^* := \big(A_{ji}^*\big)_{ij}$, acting on $H \oplus H$ by $X(\xi,\eta) := (A_{11}\xi+A_{12}\eta,\ A_{21}\xi+A_{22}\eta)$ for $\xi,\eta \in H$. A net $X_\lambda = (A_{ij}^\lambda)$ converges to $X=(A_{ij})$ in the weak operator topology on $H \oplus H$ if and only if $A_{ij}^\lambda \to A_{ij}$ weakly for every $i,j$, since the vectors $(\xi,0)$ and $(0,\xi)$, $\xi \in H$, already separate the four entries. As $M$ is weakly closed, so is $N$, and $N$ is therefore a von Neumann algebra on $H \oplus H$.

\begin{definition}[Balanced state]\label{def:balanced-state}
	Define $\Phi : N \to \mathbb{C}$ by
	\[
	\Phi(X) := \tfrac{1}{2}\big[\varphi(A_{11}) + \psi(A_{22})\big], \qquad X = (A_{ij}) \in N.
	\]
\end{definition}

$\Phi$ is a state on $N$. Linearity is immediate. For positivity, if $X = Y^*Y \geq 0$ with $Y=(B_{ij}) \in N$, then $A_{11} = B_{11}^*B_{11}+B_{21}^*B_{21} \geq 0$ and $A_{22}=B_{12}^*B_{12}+B_{22}^*B_{22}\geq 0$ in $M$, since each summand is of the form $B^*B$, so $\Phi(X) \geq 0$. Normalization holds since $\Phi(1_N) = \tfrac12[\varphi(1)+\psi(1)]=1$. $\Phi$ is faithful: writing $X^*X = \big(\sum_k A_{ki}^*A_{kj}\big)_{ij}$, $\Phi(X^*X) = \tfrac12\big[\varphi(A_{11}^*A_{11}+A_{21}^*A_{21}) + \psi(A_{12}^*A_{12}+A_{22}^*A_{22})\big]$, a sum of four nonnegative terms, so $\Phi(X^*X)=0$ forces each to vanish, and faithfulness of $\varphi,\psi$ then forces $A_{11}=A_{21}=A_{12}=A_{22}=0$, that is, $X=0$. Finally $\Phi$ is normal, being a finite linear combination of the normal functionals $\varphi,\psi$ precomposed with the weakly continuous maps $X \mapsto A_{11}$, $X \mapsto A_{22}$.

Since $\Phi$ is a normal state on the already concrete von Neumann algebra $N$, its GNS representation $\pi_\Phi(N)$ is itself weakly closed, as per Section 4.1, so $\pi_\Phi(N)$ is a von Neumann algebra on $H_\Phi$, and the GNS vector $\Xi_\Phi$ is cyclic by construction and separating because $\Phi$ is faithful, by the argument of Lemma~\ref{lem:SKM-on-MOmega}'s proof applied with $N$ in place of $M$. Theorem~\ref{thm:TT-M} and Theorem~\ref{thm:takesaki} therefore apply to $(N,\Xi_\Phi)$, producing a modular automorphism group $\sigma^\Phi$ for which $\Phi$ is $(\sigma^\Phi,1)$-KMS. To lighten notation, $\pi_\Phi$ is left implicit from here on, and elements of $N$ are identified with their images in $\pi_\Phi(N)$.

\paragraph{Corners, matrix units, and the key invariance.}
Let $p_1 := \mathrm{diag}(1,0)$, $p_2 := \mathrm{diag}(0,1) = 1_N-p_1$, and $v := \begin{pmatrix}0&0\\1&0\end{pmatrix} \in N$, where $1$ denotes the identity of $M$. A direct matrix computation gives $v^*v = p_1$, $vv^* = p_2$, so $v$ is a partial isometry implementing a Murray-von Neumann equivalence $p_1 \sim p_2$ of the type described in Section 4.2. Let $\iota_1,\iota_2 : M \to N$ be $\iota_1(A) := \mathrm{diag}(A,0)$, $\iota_2(A) := \mathrm{diag}(0,A)$, both injective $*$-homomorphisms, with $\Phi(\iota_1(A)) = \tfrac12\varphi(A)$ and $\Phi(\iota_2(A)) = \tfrac12\psi(A)$. Direct matrix multiplication gives the two identities
\begin{equation}\label{eq:corner-identities}
	v\,\iota_1(A)\,v^* = \iota_2(A), \qquad v^*\,\iota_2(A)\,v = \iota_1(A), \qquad A \in M.
\end{equation}

\begin{lemma}\label{lem:centralizer}
	$\sigma_t^\Phi(p_1) = p_1$ for every $t \in \mathbb{R}$.
\end{lemma}

\begin{proof}
	First, $\Phi(p_1X) = \Phi(Xp_1)$ for every $X=(A_{ij}) \in N$: both $p_1X$ and $Xp_1$ have $(1,1)$-entry $A_{11}$ and vanish in the row or column outside the first, and $\Phi$ reads only the $(1,1)$ and $(2,2)$ entries, so $\Phi(p_1X)=\Phi(Xp_1)=\tfrac12\varphi(A_{11})$.
	
	Second, $\Phi \circ \sigma_t^\Phi = \Phi$ for every $t$. By Proposition~\ref{prop:Delta-Omega} applied to $(N,\Xi_\Phi)$, $\Delta_\Phi^{it}\Xi_\Phi = \Xi_\Phi$ for every $t$, so
	\[
	\Phi(\sigma_t^\Phi(X)) = \big(\Xi_\Phi, \Delta_\Phi^{it}X\Delta_\Phi^{-it}\Xi_\Phi\big) = \big(\Delta_\Phi^{-it}\Xi_\Phi, X\Delta_\Phi^{-it}\Xi_\Phi\big) = (\Xi_\Phi,X\Xi_\Phi) = \Phi(X).
	\]
	
	Fix $X \in N$ and $t \in \mathbb{R}$, and apply Theorem~\ref{thm:takesaki} to the pair $(p_1,X)$: there is $F$, holomorphic and bounded on $0<\mathrm{Im}(z)<1$, continuous on the closure, with $F(t) = \Phi(p_1\sigma_t^\Phi(X))$ and $F(t+i) = \Phi(\sigma_t^\Phi(X)p_1)$. Since the first paragraph of this proof holds for every element of $N$, in particular for $\sigma_t^\Phi(X)$, $F(t) = \Phi(p_1\sigma_t^\Phi(X)) = \Phi(\sigma_t^\Phi(X)p_1) = F(t+i)$, for every $t \in \mathbb{R}$. Gluing $F$ to itself across each boundary line via this identity extends it to a bounded entire function, periodic with period $i$, which is constant by Liouville's theorem. Hence $F(t) = F(0) = \Phi(p_1X)$ for every $t \in \mathbb{R}$, that is,
	\begin{equation}\label{eq:F-constant}
		\Phi(p_1\sigma_t^\Phi(X)) = \Phi(p_1X), \qquad \forall X \in N, \ \forall t \in \mathbb{R}.
	\end{equation}
	
	Applying the second paragraph with $\sigma_{-t}^\Phi$ to the element $p_1\sigma_t^\Phi(X)$, and multiplicativity of $\sigma_{-t}^\Phi$,
	\[
	\Phi\big(p_1\sigma_t^\Phi(X)\big) = \Phi\big(\sigma_{-t}^\Phi(p_1\sigma_t^\Phi(X))\big) = \Phi\big(\sigma_{-t}^\Phi(p_1)\,X\big).
	\]
	Combined with \eqref{eq:F-constant}, $\Phi\big(\sigma_{-t}^\Phi(p_1)X\big) = \Phi(p_1X)$ for every $X \in N$. Set $Y_t := \sigma_{-t}^\Phi(p_1) - p_1$, selfadjoint since $\sigma_{-t}^\Phi$ is a $*$-automorphism applied to the selfadjoint $p_1$. The displayed identity gives $\Phi(Y_tX)=0$ for every $X \in N$, in particular for $X=Y_t$, so $\Phi(Y_t^*Y_t) = \Phi(Y_t^2) = 0$, and faithfulness of $\Phi$ forces $Y_t=0$, that is, $\sigma_{-t}^\Phi(p_1)=p_1$. As $t$ ranges over $\mathbb{R}$ this is the claim.
\end{proof}

Since $p_1$ and $p_2=1-p_1$ are both fixed by $\sigma_t^\Phi$, and $\iota_1(A) = p_1\iota_1(A)p_1$ lies entirely in the $p_1$-corner, multiplicativity of $\sigma_t^\Phi$ gives $\sigma_t^\Phi(\iota_1(A)) = p_1\sigma_t^\Phi(\iota_1(A))p_1 \in p_1Np_1 = \iota_1(M)$, so $\sigma_t^\Phi$ restricts to a well-defined automorphism of $\iota_1(M)$, and likewise of $\iota_2(M)$. Write
\[
\tau_t(A) := \iota_1^{-1}\big(\sigma_t^\Phi(\iota_1(A))\big), \qquad \rho_t(A) := \iota_2^{-1}\big(\sigma_t^\Phi(\iota_2(A))\big), \qquad A \in M.
\]

\begin{proposition}\label{prop:corner-modular}
	$\tau_t = \sigma_t^\varphi$ and $\rho_t = \sigma_t^\psi$ for every $t \in \mathbb{R}$.
\end{proposition}

\begin{proof}
	Fix $A,B \in M$ and apply Theorem~\ref{thm:takesaki} to $(\iota_1(A),\iota_1(B)) \in N$: there is $G$, holomorphic and bounded on $0<\mathrm{Im}(z)<1$, continuous on the closure, with
	\[
	G(t) = \Phi\big(\iota_1(A)\,\sigma_t^\Phi(\iota_1(B))\big) = \Phi\big(\iota_1(A\tau_t(B))\big) = \tfrac12\varphi(A\tau_t(B)),
	\]
	\[
	G(t+i) = \Phi\big(\sigma_t^\Phi(\iota_1(B))\,\iota_1(A)\big) = \Phi\big(\iota_1(\tau_t(B)A)\big) = \tfrac12\varphi(\tau_t(B)A).
	\]
	Setting $F:=2G$ exhibits $F$ as a function witnessing that $\varphi$ is $(\tau,1)$-KMS. Since $A,B$ were arbitrary, $\varphi$ is $(\tau,1)$-KMS, and the uniqueness clause of Theorem~\ref{thm:takesaki} forces $\tau_t = \sigma_t^\varphi$. The identical argument, with $\iota_2,\psi$ in place of $\iota_1,\varphi$, gives $\rho_t=\sigma_t^\psi$.
\end{proof}

\paragraph{The cocycle.}
\begin{theorem}[Connes' cocycle Radon-Nikodym theorem]\label{thm:connes-cocycle}
	Let $M$ be a von Neumann algebra, $\varphi,\psi$ faithful normal states on $M$. There is a unique family of unitaries $(D\varphi:D\psi)_t \in M$, $t \in \mathbb{R}$, such that
	\begin{enumerate}
		\item[(i)] $(D\varphi:D\psi)_{t+s} = (D\varphi:D\psi)_t\,\sigma_t^\psi\big((D\varphi:D\psi)_s\big)$ for all $s,t \in \mathbb{R}$,
		\item[(ii)] $\sigma_t^\varphi(A) = (D\varphi:D\psi)_t\,\sigma_t^\psi(A)\,(D\varphi:D\psi)_t^*$ for all $A \in M$, $t \in \mathbb{R}$.
	\end{enumerate}
	Explicitly, $(D\varphi:D\psi)_t := \iota_2^{-1}\big(v\,\sigma_t^\Phi(v)^*\big)$.
\end{theorem}

\begin{proof}
	Write $u_t := v\sigma_t^\Phi(v)^*$. Since $\sigma_t^\Phi$ is a $*$-automorphism fixing $p_1,p_2$ by Lemma~\ref{lem:centralizer}, $\sigma_t^\Phi(v)^*\sigma_t^\Phi(v) = \sigma_t^\Phi(v^*v) = \sigma_t^\Phi(p_1) = p_1$ and $\sigma_t^\Phi(v)\sigma_t^\Phi(v)^* = \sigma_t^\Phi(p_2) = p_2$, so $\sigma_t^\Phi(v)$ is a partial isometry with the same initial and final projections as $v$. Then
	\[
	u_t^*u_t = \sigma_t^\Phi(v)v^*v\sigma_t^\Phi(v)^* = \sigma_t^\Phi(v)p_1\sigma_t^\Phi(v)^* = \big(\sigma_t^\Phi(v)\sigma_t^\Phi(v)^*\big)^2 = p_2,
	\]
	using $p_1 = \sigma_t^\Phi(v)^*\sigma_t^\Phi(v)$ substituted and regrouped, and similarly $u_tu_t^* = v\sigma_t^\Phi(v)^*\sigma_t^\Phi(v)v^* = vp_1v^* = (vv^*)^2=p_2$. So $u_t$ is a unitary in $p_2Np_2 = \iota_2(M)$, and $(D\varphi:D\psi)_t := \iota_2^{-1}(u_t)$ is a well-defined unitary in $M$.
	
	For (ii), fix $A \in M$ and compute, using $\iota_2 \circ \sigma_t^\psi = \sigma_t^\Phi \circ \iota_2$, which is Proposition~\ref{prop:corner-modular},
	\[
	\iota_2\big((D\varphi:D\psi)_t\,\sigma_t^\psi(A)\,(D\varphi:D\psi)_t^*\big) = u_t\,\sigma_t^\Phi(\iota_2(A))\,u_t^* = v\sigma_t^\Phi(v)^*\sigma_t^\Phi(\iota_2(A))\sigma_t^\Phi(v)v^*.
	\]
	By \eqref{eq:corner-identities}, $v^*\iota_2(A)v = \iota_1(A)$, so $\sigma_t^\Phi(v)^*\sigma_t^\Phi(\iota_2(A))\sigma_t^\Phi(v) = \sigma_t^\Phi\big(v^*\iota_2(A)v\big) = \sigma_t^\Phi(\iota_1(A)) = \iota_1(\sigma_t^\varphi(A))$, the last equality by Proposition~\ref{prop:corner-modular}. Substituting and applying \eqref{eq:corner-identities} once more, $v\,\iota_1(\sigma_t^\varphi(A))\,v^* = \iota_2(\sigma_t^\varphi(A))$. Hence $\iota_2\big((D\varphi:D\psi)_t\sigma_t^\psi(A)(D\varphi:D\psi)_t^*\big) = \iota_2(\sigma_t^\varphi(A))$, and injectivity of $\iota_2$ gives (ii).
	
	For (i), using $u_t^*u_t = p_1$, the partial isometry identity $vp_1=v$, and $\sigma_t^\Phi(\sigma_s^\Phi(v)^*) = \sigma_t^\Phi(\sigma_s^\Phi(v^*))$,
	\[
	u_t\sigma_t^\Phi(u_s) = v\sigma_t^\Phi(v)^* \cdot \sigma_t^\Phi\big(v\sigma_s^\Phi(v)^*\big) = v\big[\sigma_t^\Phi(v)^*\sigma_t^\Phi(v)\big]\sigma_t^\Phi(\sigma_s^\Phi(v^*)) = (vp_1)\,\sigma_{t+s}^\Phi(v^*) = v\,\sigma_{t+s}^\Phi(v^*) = u_{t+s}.
	\]
	Since $\iota_2\circ\sigma_t^\psi=\sigma_t^\Phi\circ\iota_2$, applying $\iota_2$ to (i)'s right-hand side gives $\iota_2\big((D\varphi:D\psi)_t\sigma_t^\psi((D\varphi:D\psi)_s)\big) = u_t\sigma_t^\Phi(u_s) = u_{t+s} = \iota_2\big((D\varphi:D\psi)_{t+s}\big)$, and injectivity of $\iota_2$ gives (i).
	
	Uniqueness is proved in \cite[TODO: collocazione precisa]{TODO-Connes73}, in the equivalent formulation of \cite[TODO: collocazione precisa]{TODO-Stratila}.
\end{proof}

\begin{corollary}[Chain rule]\label{cor:chain-rule}
	For faithful normal states $\varphi,\psi,\chi$ on $M$, $(D\varphi:D\chi)_t = (D\varphi:D\psi)_t\,(D\psi:D\chi)_t$ for every $t \in \mathbb{R}$.
\end{corollary}

\begin{proof}
	Write $u_t := (D\varphi:D\psi)_t$, $v_t:=(D\psi:D\chi)_t$, $w_t := u_tv_t$. By property (ii) applied twice, for $A \in M$,
	\[
	\sigma_t^\varphi(A) = u_t\sigma_t^\psi(A)u_t^* = u_t\big(v_t\sigma_t^\chi(A)v_t^*\big)u_t^* = w_t\,\sigma_t^\chi(A)\,w_t^*,
	\]
	so $w_t$ satisfies property (ii) for the pair $(\varphi,\chi)$. By property (i) applied to each factor and property (ii) for $(\psi,\chi)$ applied to the element $u_s$,
	\[
	w_{t+s} = u_{t+s}v_{t+s} = u_t\sigma_t^\psi(u_s)\cdot v_t\sigma_t^\chi(v_s) = u_t\big[\sigma_t^\psi(u_s)v_t\big]\sigma_t^\chi(v_s) = u_t\big[v_t\sigma_t^\chi(u_s)\big]\sigma_t^\chi(v_s) = w_t\,\sigma_t^\chi(u_sv_s) = w_t\sigma_t^\chi(w_s),
	\]
	so $w_t$ also satisfies property (i) for $\sigma^\chi$. By the uniqueness clause of Theorem~\ref{thm:connes-cocycle}, $w_t = (D\varphi:D\chi)_t$.
\end{proof}

\begin{theorem}\label{thm:out-independence}
	For every pair of faithful normal states $\varphi,\psi$ on a von Neumann algebra $M$ and every $t \in \mathbb{R}$, $\sigma_t^\varphi$ and $\sigma_t^\psi$ represent the same element of $\mathrm{Out}(M)$.
\end{theorem}

\begin{proof}
	By property (ii) of Theorem~\ref{thm:connes-cocycle}, $\sigma_t^\varphi = \mathrm{Ad}\big((D\varphi:D\psi)_t\big) \circ \sigma_t^\psi$, as per Definition~\ref{def:inn-out}, since $(D\varphi:D\psi)_t$ is a unitary in $M$. Hence $\sigma_t^\varphi$ and $\sigma_t^\psi$ differ by an inner automorphism, and coincide in $\mathrm{Aut}(M)/\mathrm{Inn}(M)$.
\end{proof}

\begin{remark}
	Theorem~\ref{thm:out-independence} is exactly the well-posedness required by Definition~\ref{def:thermal-time}: whichever faithful normal state a system described by $M$ is found in, the modular automorphism group it determines is, as a class in $\mathrm{Out}(M)$, the same physical flow. The thermal time hypothesis identifies this class with time itself. Section 4.9 refines this picture further, showing that the modular spectrum, an invariant of $M$ alone, already constrains which such flows are possible before any state is chosen at all.
\end{remark}

\section{Reformulating the Axiomatic Framework of Section 3.7 via the Modular Flow}\label{sec:axioms-modular}

Definition 3.7.2 packaged the geometric reconstruction of quantum mechanics into four axioms, with axiom (ii) requiring time evolution on each sector $PH_\alpha$ to be the Hamiltonian flow of the mean value function of a distinguished selfadjoint element $H \in A_{sa}$. Nothing in that axiom, or in the theorems of Chapter 3 leading up to it, explains where $H$ comes from. It is simply postulated, exactly as a Hamiltonian is postulated in the Dirac-von Neumann formalism itself, so the geometric axiomatization gains nothing on this one point beyond trading a postulate for an equivalent postulate phrased geometrically. For a system described by a type III von Neumann algebra, per Section 4.2, the gap is more than cosmetic, since there is no trace with which to write $H$ down as a density-matrix logarithm in the first place. This section shows that the modular automorphism group of Section 4.4, attached to $M$ by nothing more than a single faithful normal state and no further data, supplies a generator satisfying exactly the geometric requirement axiom (ii) imposes, Kähler flow of an observable function on $PH_\omega$, derived rather than postulated. Section 4.6 already shows, however, that this generator need not be the original Hamiltonian itself, a point this section makes precise before proposing the replacement axiom.

\begin{definition}[Modular Hamiltonian]\label{def:modular-hamiltonian}
	Let $M$ be a von Neumann algebra, $\omega$ a faithful normal state on $M$, $(H_\omega,\pi_\omega,\Omega_\omega)$ its GNS representation, as per Theorem 2.7, and $\Delta_\omega$ the modular operator of $\pi_\omega(M)$ relative to $\Omega_\omega$, as per Definition~\ref{def:modular-M}, applied with $\pi_\omega(M)$ in the role of $M$. Section 4.1 already identifies $\pi_\omega(M)$ as a von Neumann algebra on $H_\omega$, so this specialization is legitimate. By Theorem 2.10, applied to the strongly continuous one-parameter unitary group $(\Delta_\omega^{it})_{t \in \mathbb{R}}$, there is a unique selfadjoint operator $L_\omega$ on $H_\omega$ such that $\Delta_\omega^{it} = e^{itL_\omega}$ for every $t \in \mathbb{R}$. Fix $\beta > 0$. The modular Hamiltonian of $(M,\omega)$ at inverse temperature $\beta$ is $h_\omega := -\beta^{-1}L_\omega$.
\end{definition}

\begin{proposition}[Modular flow as a Kähler isomorphism group]\label{prop:modular-kahler}
	Let $M$, $\omega$, $\Delta_\omega$, $h_\omega$ be as per Definition~\ref{def:modular-hamiltonian}, and let $PH_\omega$ be the projective space of $H_\omega$, as per Definition 3.2.2, carrying the Fubini-Study triad $(J,g,\omega)$ of Equation (3.24).
	\begin{enumerate}
		\item[(a)] For every $t \in \mathbb{R}$, $[x] \mapsto [\Delta_\omega^{it}x]$ is a Kähler isomorphism of $PH_\omega$, and $t \mapsto \Delta_\omega^{it}$ is a one-parameter group of such isomorphisms, generated in the sense of part 2 of Proposition 3.3.6 by $H := \beta\hbar h_\omega$, that is, $\Delta_\omega^{it}=e^{-itH/\hbar}$.
		\item[(b)] If $h_\omega \in B(H_\omega)$, the group of part (a) is the Hamiltonian flow, with respect to $g$, of the observable function $\langle H \rangle \in K(PH_\omega,\mathbb{R})$, as per Definitions 3.3.1 and 3.5.1, and its generating vector field is Killing for $g$, as per Definition 3.5.2.
	\end{enumerate}
\end{proposition}

\begin{proof}
	(a) By Proposition~\ref{prop:polar-SK}, $\Delta_\omega$ is positive and selfadjoint, so $\Delta_\omega^{it}$ is unitary on $H_\omega$ for every $t \in \mathbb{R}$. By part 1 of Proposition 3.3.6, every unitary operator on $H_\omega$ induces a Kähler isomorphism of $PH_\omega$, so $[x]\mapsto[\Delta_\omega^{it}x]$ is one for every fixed $t$, and since $\Delta_\omega^{it}\Delta_\omega^{is}=\Delta_\omega^{i(t+s)}$, the projectivized maps form a one-parameter group. By Definition~\ref{def:modular-hamiltonian}, $\Delta_\omega^{it}=e^{itL_\omega}=e^{-i\beta t h_\omega}$, which equals $e^{-itH/\hbar}$ for $H:=\beta\hbar h_\omega$, matching the form of part 2 of Proposition 3.3.6.
	
	(b) If $h_\omega$ is bounded, so is $H=\beta\hbar h_\omega \in B(H_\omega)$, and $H$ is selfadjoint since $h_\omega$ is. Proposition 3.3.2 then applies verbatim to $H$, identifying the Hamiltonian flow of $\langle H\rangle$ with $[x]\mapsto[e^{-itH/\hbar}x]=[\Delta_\omega^{it}x]$, the flow of part (a). Since this flow preserves $g$ by part (a), being Kähler and not merely symplectic, its generating vector field is Killing for $g$ by Definition 3.5.2, consistently with Theorem 3.9, which independently guarantees that an observable function such as $\langle H\rangle$ could only have generated a Killing vector field.
\end{proof}

Part (b) is the precise sense in which the modular flow recovers axiom (ii) of Definition 3.7.2, once the hypothesis $h_\omega \in B(H_\omega)$ holds. The next remark checks this hypothesis in the one setting where $h_\omega$ has already been computed, and shows that the recovered generator, while bounded, is not simply the original Hamiltonian read back.

\begin{remark}
	In the setting of Section 4.6, $M=M_n(\mathbb{C})$, $\omega=\omega_\beta$, and $\Delta_{\omega_\beta}^{it}(X)=\rho^{it}X\rho^{-it}$ for $\rho=e^{-\beta H}$, the original Hamiltonian $H$ of that example. Differentiating at $t=0$ gives $L_{\omega_\beta}(X) = -\beta(HX-XH)$, so $h_{\omega_\beta}(X) = HX-XH$, the commutator with $H$, a bounded operator on $H_{\omega_\beta}$ since $H$ is bounded. The selfadjoint element recovered by Proposition~\ref{prop:modular-kahler}(b) is therefore $\beta\hbar h_{\omega_\beta} = \beta\hbar(HX-XH)$, which is not $\pi_{\omega_\beta}(H)=L_H$, left multiplication by $H$ alone, but a combination of $L_H$ and $-R_H$, mixing $M=\pi_{\omega_\beta}(M_n(\mathbb{C}))$ and its commutant $M'$ exactly as Section 4.4 anticipated when it observed that the modular objects of $K_M$ and $K_{M'}$ are not independent. Axiom (ii) of Definition 3.7.2 only requires $H \in B(H_\omega)$ selfadjoint, not $H \in \pi_\omega(M)$, so this mismatch does not violate the axiom, but it does mean the two notions of dynamics, ordinary Heisenberg evolution under $H$ and modular flow at the state $\omega_\beta$, are genuinely different flows on $PH_{\omega_\beta}$, agreeing only through the relation that Proposition~\ref{prop:modular-kahler} makes precise.
\end{remark}

With this distinction in view, axiom (ii) of Definition 3.7.2 can now be replaced outright by a version that no longer presupposes a distinguished Hamiltonian as separate data.

\begin{definition}[Modular dynamics]\label{def:modular-dynamics-axiom}
	An algebraic-geometric quantum system $(A,S,H)$, as per Definition 3.7.2, satisfies modular dynamics if axiom (ii) is replaced by the following. On each superselection sector $PH_\alpha$, fixed by a faithful normal state $\omega$ on $A$ restricted to that sector, time evolution is the projectivization of the modular automorphism group $\Delta_\omega^{it}$ of Definition~\ref{def:modular-automorphism}, for a fixed choice of $\beta>0$.
\end{definition}

\begin{remark}
	Definition~\ref{def:modular-dynamics-axiom} requires only the algebraic core, axiom (i) of Definition 3.7.2, and the choice of a faithful normal state on each sector, already needed for axiom (iii) to make operational sense, as per Remark 3.7.1. No selfadjoint element of $A_{sa}$ need be singled out, and the construction is available whether or not $A$ carries a trace, in contrast with Definition 3.7.2 as originally stated, which is silent on where $H$ comes from in a type III sector. By Proposition~\ref{prop:modular-kahler}(a), modular dynamics always produces a well-defined one-parameter group of Kähler isomorphisms of $PH_\omega$, so Definition~\ref{def:modular-dynamics-axiom} is a legitimate strengthening of axiom (i) alone, whether or not axiom (ii) as originally stated can also be satisfied. By part (b), whenever a bounded modular Hamiltonian exists, modular dynamics coincides with the original axiom (ii) for the specific generator $H=\beta\hbar h_\omega$, so Definition~\ref{def:modular-dynamics-axiom} subsumes rather than contradicts Definition 3.7.2 in every case the latter already covered. Theorem~\ref{thm:out-independence} guarantees that this entire construction depends on the choice of $\omega$ among faithful normal states only up to an inner automorphism of $A$, so modular dynamics is, in the sense of Section 4.7, a single physical fact about the sector rather than an artifact of which state happened to fix it.
\end{remark}

\section{The Modular Spectrum and the Fine Classification of Type III}\label{sec:modular-spectrum}

Section 4.2 classified factors into types I, II, and III by the presence or absence of a trace, and left every type III factor on the same footing, distinguished from type I and II but not from one another. This is a genuine gap, not a mere labeling convenience, since type III factors arising from different physical systems, for instance from different local algebras in quantum field theory, are not isomorphic to each other despite sharing the single qualitative feature of having no trace. This section shows that the modular operator, already exhibited in Sections 4.4 through 4.6 as a substitute for the dynamics a trace would otherwise supply, carries in addition a spectral invariant fine enough to split type III into three genuinely different classes, closing the gap left open in Section 4.2.

\begin{lemma}\label{lem:one-in-spectrum}
	Let $M$ be a von Neumann algebra and $\omega$ a faithful normal state on $M$, with modular operator $\Delta_\omega$ relative to the GNS vector $\Omega_\omega$, as per Definition~\ref{def:modular-M} applied to $\pi_\omega(M)$. Then $1 \in \mathrm{Spec}(\Delta_\omega)$.
\end{lemma}

\begin{proof}
	By Proposition~\ref{prop:Delta-Omega}, $\Delta_\omega\Omega_\omega = \Omega_\omega$, and $\Omega_\omega \neq 0$ since it has unit norm as per Theorem 2.7. Hence $1$ is an eigenvalue of $\Delta_\omega$, and every eigenvalue lies in the spectrum.
\end{proof}

\begin{definition}[Connes' $S$-invariant]\label{def:connes-S}
	Let $M$ be a factor, as per Section 4.2, and let $\mathcal{N}(M)$ denote the set of faithful normal states on $M$. The $S$-invariant of $M$ is
	\[
	S(M) := \bigcap_{\omega \in \mathcal{N}(M)} \mathrm{Spec}(\Delta_\omega) \ \subseteq \ [0,\infty).
	\]
\end{definition}

\noindent By Lemma~\ref{lem:one-in-spectrum}, $1 \in S(M)$ for every factor $M$, so $S(M)$ is never empty.

The classification itself is Connes' theorem, obtained by a substantially deeper analysis of the modular automorphism group than Section 4.7 develops, using the Connes spectrum of a one-parameter automorphism group rather than the operator spectrum of a single $\Delta_\omega$ to show that $S(M)$, defined above as an intersection over every state, is already realized by any single faithful normal state once $M$ is a factor. It is quoted here as a black box, with the technical hypothesis of $M$ having separable predual, equivalently countable decomposability, made explicit since every construction in this chapter has in fact produced such an $M$, being built from a single GNS representation at each stage.

\begin{theorem}[Connes, 1973]\label{thm:connes-classification}
	Let $M$ be a factor with separable predual.
	\begin{enumerate}
		\item[(a)] $M$ is semifinite, that is, of type I or II as per Section 4.2, if and only if $S(M) = \{1\}$.
		\item[(b)] If $M$ is of type III, then $S(M)\setminus\{0\}$ is a closed subgroup of the multiplicative group $(0,\infty)$, and is accordingly one of exactly three possibilities:
		\begin{itemize}
			\item $S(M)\setminus\{0\} = \{1\}$, in which case $M$ is said to be of type $\mathrm{III}_0$,
			\item $S(M)\setminus\{0\} = \{\lambda^n \mid n \in \mathbb{Z}\}$ for a unique $\lambda \in (0,1)$, in which case $M$ is said to be of type $\mathrm{III}_\lambda$,
			\item $S(M)\setminus\{0\} = (0,\infty)$, in which case $M$ is said to be of type $\mathrm{III}_1$.
		\end{itemize}
	\end{enumerate}
\end{theorem}
\noindent This is proved in \cite[TODO: collocazione precisa]{TODO-Connes73}, with a textbook treatment in \cite[TODO: collocazione precisa]{TODO-Stratila}.

\begin{remark}
	Theorem~\ref{thm:connes-classification} is the precise sense in which type III, coarsely defined in Section 4.2 as the absence of a trace, is not a single phenomenon. A type $\mathrm{III}_0$ factor has a trivial modular spectrum away from $0$, the closest a type III factor can come to being semifinite without a trace actually existing, while a type $\mathrm{III}_1$ factor realizes every positive real as a modular eigenvalue ratio, the opposite extreme. The intermediate family $\mathrm{III}_\lambda$, $\lambda \in (0,1)$, is discrete, parametrized by a single modulus, and is exactly the case realized by the Gibbs states of Section 4.6 once the finite-dimensional setting of that example is replaced by an infinite one, as noted again below.
\end{remark}

That Theorem~\ref{thm:connes-classification} manages to extract, from an intersection over every faithful normal state, an invariant already visible in any single one of them, is consistent with Theorem~\ref{thm:out-independence}: the modular automorphism groups of two different states differ only by conjugation with a cocycle of unitaries already in $M$, so no single choice of $\omega$ can see more of the algebra $M$ through $\Delta_\omega$ than any other choice does. Making this heuristic into the actual proof of Theorem~\ref{thm:connes-classification}, however, requires tracking how the Connes cocycle of Section 4.7 acts on spectral projections of $\Delta_\omega$ itself, a refinement of the cocycle theorem not developed here.

\begin{remark}[The semifinite case verified directly]
	Section 4.2 identified $M_n(\mathbb{C})$ as a semifinite, indeed type I, factor, and Theorem~\ref{thm:connes-classification}(a) predicts $S\big(M_n(\mathbb{C})\big) = \{1\}$. This is checked directly from Section 4.6 with no appeal to Theorem~\ref{thm:connes-classification} itself. By Lemma~\ref{lem:one-in-spectrum}, $1 \in S\big(M_n(\mathbb{C})\big)$. Conversely, the tracial state $\omega_0$ of the Remark closing Section 4.6 is a faithful normal state on $M_n(\mathbb{C})$ with $\Delta_{\omega_0} = 1$, the identity operator, so $\mathrm{Spec}(\Delta_{\omega_0}) = \{1\}$, and since $S\big(M_n(\mathbb{C})\big)$ is an intersection including this term,
	\[
	S\big(M_n(\mathbb{C})\big) \ \subseteq \ \mathrm{Spec}(\Delta_{\omega_0}) = \{1\}.
	\]
	Combined with the reverse inclusion, $S\big(M_n(\mathbb{C})\big) = \{1\}$, confirming Theorem~\ref{thm:connes-classification}(a) in this example without circularity.
\end{remark}

\begin{remark}[Type $\mathrm{III}_\lambda$ anticipated]
	The computation of Section 4.6 identified $\mathrm{Spec}(\Delta_{\omega_\beta})$, for the Gibbs state on $M_n(\mathbb{C})$ at inverse temperature $\beta$, with the finite set $\{\mu_p/\mu_q \mid p,q=1,\dots,n\}$ of ratios of eigenvalues $\mu_k = e^{-\beta\lambda_k}/Z_\beta$ of the density matrix. On the finite-dimensional algebra $M_n(\mathbb{C})$ this set is always finite, never a nontrivial closed subgroup of $(0,\infty)$, consistently with $M_n(\mathbb{C})$ being semifinite rather than type III. Were the same construction carried out on an infinite tensor product of matrix algebras with a fixed ratio of level spacings repeating at every site, an infinite-dimensional algebra outside the scope of Section 4.6, the analogous set of ratios would already be the discrete group $\{\lambda^n \mid n \in \mathbb{Z}\}$ for $\lambda$ fixed by the spacing, and the resulting hyperfinite factor is the standard example of type $\mathrm{III}_\lambda$ quoted alongside Theorem~\ref{thm:connes-classification} in \cite[TODO: collocazione precisa]{TODO-Stratila}.
\end{remark}

\section{Purification Revisited: Thermofield Double and Stratified State Spaces}\label{sec:purification-revisited}

Remark 3.7.7 closed Chapter 3 by observing that the space of density operators on a fixed sector is not itself a smooth manifold, its local dimension jumping with the rank of the state, and that the Gelfand-Naimark-Segal construction of Theorem 2.7 sidesteps this obstruction by representing every state, pure or mixed, as a vector state of some Hilbert space $H_\omega$, landing on the honest Kähler manifold $PH_\omega$ of Theorem 3.4 regardless of the rank of the state purified. Chapter 4 has used this device throughout without remarking on it: the Gibbs state $\omega_\beta$ of Section 4.6 is mixed for every finite $\beta$, yet its GNS vector $\Omega_\beta$ was handled exactly as a cyclic and separating vector for a von Neumann algebra, with no distinction drawn between this case and a pure state. This section makes the resulting purification, widely known in physics as the thermofield double construction, precise within the framework already built, and asks what the modular operator of this chapter adds to the purely topological resolution Remark 3.7.7 already supplied.

\begin{proposition}[Purification is always pure]\label{prop:purification-pure}
	Let $M$ be a von Neumann algebra on $H$, $\omega$ a faithful normal state on $M$, and $(H_\omega,\pi_\omega,\Omega_\omega)$ its GNS representation, as per Theorem 2.7. The vector state $\tilde\omega(T) := (\Omega_\omega,T\Omega_\omega)$, $T \in B_0(H_\omega)$, is pure, and $[\Omega_\omega] \in PH_\omega$ is a well-defined point of the projective Hilbert space of Definition 3.2.2, regardless of whether $\omega$ itself is pure on $M$.
\end{proposition}

\begin{proof}
	By Proposition 3.2.4, $PH_\omega$ is identified with the pure state space of $B_0(H_\omega)$, under the correspondence sending a unit vector $\Psi \in H_\omega$, up to phase, to the pure vector state $T \mapsto (\Psi,T\Psi)$. Since $\Omega_\omega$ is a unit vector by Theorem 2.7, this applies directly to $\Omega_\omega$, regardless of the algebra $M$ or the state $\omega$ that produced it as a GNS vector.
\end{proof}

Proposition~\ref{prop:purification-pure} is the general statement behind the explicit construction of Section 4.6: $\omega_\beta$ is a mixed state of $M_n(\mathbb{C})$ for every finite $\beta$, yet $\Omega_\beta$ is pure as a state of $B_0(H_{\omega_\beta})$, and in fact of $M \vee M'$ there, since left and right multiplications together act irreducibly on $H_{\omega_\beta} = M_n(\mathbb{C})$. The thermofield double construction is exactly this passage, from a mixed state $\omega$ of $M$ to the pure state $\tilde\omega$ of the doubled system $B_0(H_\omega) \supseteq M \vee M'$ that reproduces $\omega$ on $M$ upon restriction, with $M'$ playing the role of a purifying ancilla. How far this purification is from being symmetric between $M$ and $M'$ is exactly what the modular operator measures.

\begin{proposition}[Modular triviality and traciality]\label{prop:Delta-trivial-tracial}
	Let $M$ be a von Neumann algebra with cyclic and separating vector $\Omega$, and let $\omega(\cdot) := (\Omega,\cdot\,\Omega)$ be the state it induces. Let $\Delta$ be the modular operator of $M$ relative to $\Omega$, as per Definition~\ref{def:modular-M}. Then $\Delta = 1$ if and only if $\omega$ is a trace on $M$, that is, $\omega(AB) = \omega(BA)$ for every $A,B \in M$.
\end{proposition}

\begin{proof}
	For $A,B \in M$, write $\xi_A := A\Omega+A^*\Omega$, $\xi_B := B\Omega+B^*\Omega$, both in $K_M$ as per Definition~\ref{def:KM-subspace}. The computation of Lemma~\ref{lem:MMp-orthogonal}, carried out with $A,B \in M$ in place of $A \in M, A' \in M'$, no longer has access to commutativity of $A$ and $B$, and instead gives
	\begin{equation}\label{eq:xiA-xiB}
		(\xi_A,\xi_B) = \omega\big((A+A^*)(B+B^*)\big).
	\end{equation}
	Every selfadjoint element $C \in M$ is of this form, $C=A+A^*$ for $A:=C/2$, so \eqref{eq:xiA-xiB} says that $\mathrm{Im}(\xi_A,\xi_B)=0$ for every $A,B\in M$ if and only if $\omega(C_1C_2)=\omega(C_2C_1)$ for every selfadjoint $C_1,C_2 \in M$, since $\mathrm{Im}(\xi_A,\xi_B)=0$ is equivalent to $(\xi_A,\xi_B)=(\xi_B,\xi_A)$. By the Cartesian decomposition of Corollary 2.1.28, applied to a general pair $X,Y\in M$ and expanded bilinearly against this selfadjoint-pair identity, the latter holds for every selfadjoint pair if and only if $\omega(XY)=\omega(YX)$ for every $X,Y \in M$, that is, if and only if $\omega$ is a trace.
	
	It remains to relate $\mathrm{Im}(\xi,\eta)=0$ for every $\xi,\eta \in K_M$ to $\Delta=1$. For $\xi,\eta \in H$, expanding the inner product of $\xi+i\eta$ and of $\xi-i\eta$ with themselves gives
	\[
	\|\xi+i\eta\|^2 - \|\xi-i\eta\|^2 = -4\,\mathrm{Im}(\xi,\eta),
	\]
	so, by Definition~\ref{def:tomita-K}, $S_{K_M}(\xi+i\eta)=\xi-i\eta$ is norm-preserving on $D(S_{K_M})=K_M+iK_M$ if and only if $\mathrm{Im}(\xi,\eta)=0$ for every $\xi,\eta\in K_M$, and this holds for every $\xi,\eta \in K_M$, the real-linear span of the generators $\xi_A$, exactly when it holds for every pair of generators, by bilinearity and continuity of $\mathrm{Im}(\cdot,\cdot)$.
	
	Suppose $S_{K_M}$ is isometric on $K_M+iK_M$. If $(\zeta_n)\subset K_M+iK_M$ is Cauchy, then $(S_{K_M}\zeta_n)$ is Cauchy as well, by isometry, so $\zeta_n\to\zeta$ and $S_{K_M}\zeta_n\to\chi$ for some $\zeta,\chi\in H$, and closedness of $S_{K_M}$, as per Lemma~\ref{lem:SK-properties}, gives $\zeta\in K_M+iK_M$. Hence $K_M+iK_M$ is complete, so closed, and dense by Lemma~\ref{lem:KM-density}, so $K_M+iK_M=H$ and $S_{K_M}$ is an everywhere-defined isometric antilinear involution, that is, an antiunitary involution. Its polar decomposition is then $S_{K_M}=S_{K_M}\cdot 1$, so $J=S_{K_M}$ and $\Delta=1$ by uniqueness of the polar decomposition.
	
	Conversely, if $\Delta=1$, then $S_{K_M}=J\Delta^{1/2}=J$ is antiunitary and everywhere defined, so isometric on $K_M+iK_M\subseteq H$ in particular, giving $\mathrm{Im}(\xi,\eta)=0$ for every $\xi,\eta\in K_M$ by the displayed identity above.
	
	Chaining the two equivalences, $\Delta=1$ if and only if $\omega$ is a trace on $M$.
\end{proof}

\begin{remark}
	Proposition~\ref{prop:Delta-trivial-tracial} identifies exactly when the thermofield double construction above is as symmetric as possible between $M$ and $M'$: the maximally mixed state $\omega_0$ of Section 2.7 and the $\beta\to0$ limit of Section 4.6 both have $\Delta=1$, consistently with both being tracial, while a generic faithful normal state, such as $\omega_\beta$ for $\beta\in(0,\infty)$, has $\Delta\neq1$ precisely because it is not tracial, $\omega_\beta(AB)\neq\omega_\beta(BA)$ in general once $H$ fails to be a scalar. The modular operator, in this sense, measures how far the purification $\Omega$ of a given state is from the symmetric purification underlying the maximally entangled thermofield double, with the modular automorphism group of Section 4.4 quantifying this asymmetry as a flow rather than as a single number.
\end{remark}

\begin{remark}[Exploratory: a modular metric on mixed states]
	Proposition~\ref{prop:Delta-trivial-tracial} and Section 4.6 together suggest, without establishing, a direction left open by this thesis. Remark 3.7.7 bypasses the stratified space of density matrices on a fixed sector by purifying each mixed state into a pure point of the honest Kähler manifold $PH_\omega$, but different mixed states $\omega,\omega'$ on the same $M$ generally purify into different Hilbert spaces $H_\omega,H_{\omega'}$, not into points of one common manifold, so the Fubini-Study geometry of Chapter 3 says nothing directly about distances between mixed states themselves. The modular theory of this chapter attaches to every such $\omega$ an operator $\Delta_\omega$ on its own $H_\omega$, and Theorem~\ref{thm:connes-cocycle} already supplies a unitary comparing the GNS representations of two states $\varphi,\psi$ on the same $M$ at every modular time $t$. Whether this cocycle, or the relative modular operator obtained by polarizing the construction of Section 4.4 between two cyclic vectors rather than one, induces a genuine Riemannian or Kähler-type structure directly on the state space of $M$, restricting on pure states to the Fubini-Study metric of Section 3.2 and reducing on commuting states to some classical statistical distance, is not addressed here. The relative entropy of Araki, built from exactly this relative modular operator, is the natural candidate quantity to investigate first, \cite[TODO: collocazione precisa]{TODO-Araki}.
\end{remark}

\section{Discussion and Outlook}\label{sec:discussion-ch4}

Section 4.2 opened this chapter with an obstruction: a von Neumann algebra of type III carries no trace, so neither the density-matrix dynamics of ordinary quantum statistical mechanics nor, more fundamentally, the symmetrized construction underlying Definition 1.6.23 is available to furnish it with a privileged notion of time evolution. Every von Neumann algebra encountered in the physical literature that models an infinite system, from the local algebras of quantum field theory to black hole horizons, is in fact of type III, so this obstruction is not a pathology of some exotic edge case but the generic situation. The remainder of this chapter answered the obstruction, not by restoring a trace, which does not exist, but by replacing it with a structure that exists on every von Neumann algebra carrying a cyclic and separating vector, trace or no trace: the modular operator.

Section 4.3 built this structure at its most elementary level, with no algebra in sight at all, following Rieffel and Van Daele: a closed real subspace of a Hilbert space satisfying two purely geometric conditions already determines, through an antilinear involution and its polar decomposition, a positive operator $\Delta_K$ and an antiunitary $J_K$, with $\Delta_K^{it}$ leaving the subspace invariant and $J_K$ exchanging it with its symplectic complement. Section 4.4 specialized this to the real subspace generated by a von Neumann algebra $M$ and a cyclic separating vector $\Omega$, proving the elementary half of the resulting Tomita-Takesaki theorem in full, the density argument, the symplectic self-intersection argument, and the orthogonality computation of Lemma~\ref{lem:MMp-orthogonal}, while citing Rieffel and Van Daele for the single step this thesis could not reduce further, the promotion of invariance of the subspace $K_M$ to invariance of the algebra $M$ itself under $\Delta^{it}(\cdot)\Delta^{-it}$. Section 4.5 then justified calling the resulting flow $\sigma^\Omega$ a dynamics rather than bookkeeping, by proving it satisfies, and uniquely so, the Kubo-Martin-Schwinger equilibrium condition for the state $\omega$ induced by $\Omega$, with the converse Haag-Hugenholtz-Winnink direction quoted from the literature. Section 4.6 made every one of these objects explicit on $M_n(\mathbb{C})$, verifying by direct computation, with no citation required anywhere in that section, that the modular flow of a Gibbs state is ordinary Heisenberg evolution rescaled by the inverse temperature, and that the construction degenerates correctly to the identity exactly when the state degenerates to the trace.

Section 4.7 asked whether $\sigma^\Omega$ depends essentially on the particular vector $\Omega$ chosen to represent a state, or on the state at all, and proved, via the two-by-two balanced algebra of Connes, that it does not: the modular automorphism groups of any two faithful normal states on the same $M$ differ only by an inner automorphism, an explicit unitary cocycle constructed in full in this thesis with only its uniqueness left to the literature. This is what let Section 4.7 state the thermal time hypothesis of Connes and Rovelli as a well-posed proposal rather than an arbitrary one, and let Section 4.8 use the modular flow in place of a postulated Hamiltonian in the axiomatic framework of Section 3.7, with the honest qualification, visible only once the Gibbs example of Section 4.6 was worked out explicitly, that the recovered generator is a modular Hamiltonian built from the commutator with the original Hamiltonian, not that Hamiltonian itself. Section 4.9 closed the gap this chapter opened with: the modular spectrum $S(M)$, Connes' invariant built from the same operator $\Delta_\omega$ throughout, splits type III into the three classes $\mathrm{III}_0$, $\mathrm{III}_\lambda$, $\mathrm{III}_1$ that the coarse trace-based trichotomy of Section 4.2 could not distinguish, so the absence of a trace, the problem this chapter started from, is answered by the same construction, the modular operator, that resolves it. Section 4.10 returned to the purification machinery of Chapter 3, identifying exactly when a von Neumann algebra's modular operator trivializes, precisely when the state purified is tracial, and sketched, without developing, the question of whether the apparatus of this chapter extends the Fubini-Study geometry of Chapter 3 from pure states to a genuine metric structure on mixed states.

Three claims can be drawn from the chapter as a whole. First, the modular automorphism group is a mathematically canonical substitute for Hamiltonian time evolution, attached to any von Neumann algebra with a cyclic and separating vector by nothing more than that vector, available precisely where the classical construction of Definition 1.6.23 is not. Second, this substitute is not arbitrary: Section 4.7's cocycle theorem shows it to be a fact about the algebra and the state up to inner automorphism, not an artifact of auxiliary choices, which is the mathematical content underwriting the thermal time hypothesis. Third, the same operator that supplies this dynamics also supplies, through Connes' spectral invariant of Section 4.9, a finer classification of the algebra itself than the trace-based trichotomy with which the chapter began, so dynamics and classification turn out to be two readings of one object rather than two separate pieces of structure.

Measured against the two full proofs this chapter set out to deliver, it met one without qualification and the other only in part. The two-by-two matrix proof of Connes' cocycle theorem in Section 4.7, both the existence of $(D\varphi:D\psi)_t$ with its explicit formula and both of its defining properties, is proved from nothing beyond Sections 4.4 and 4.5, with only the separate fact of its uniqueness, needed for the chain rule, left to the literature. The Tomita-Takesaki theorem of Section 4.4 did not reach the same standard: every argument specific to the algebra $M$, the construction of $K_M$, its standardness, the identification $K_M'=K_{M'}$ up to its hard commutation-theorem half, and the agreement of $S_{K_M}$ with the naive Tomita operator on $M\Omega$, is proved in full, but the passage from subspace-level to algebra-level invariance remains a citation, despite repeated attempts in this thesis to reduce it further. A reader who wants the standard subspace route to be elementary without exception should consult Rieffel and Van Daele directly on this one point.

Beyond this, the chapter has a narrower scope than its subject matter allows. Every explicit computation, Section 4.6's Gibbs state and the semifinite verification of Section 4.9, took place on $M_n(\mathbb{C})$, a type I factor, so no type III example was ever constructed by hand, only gestured at through the infinite tensor product sketched in the second remark of Section 4.9. The exploratory remark closing Section 4.10, on whether a relative modular operator induces a genuine metric on the space of mixed states extending the Fubini-Study geometry of Section 3.2, through Araki's relative entropy or otherwise, was deliberately left open rather than pursued. Both of these point toward the same place the Conclusion of this thesis already identified as unfinished business, extending the algebraic-geometric framework to infinite-dimensional Hilbert spaces and to mixed states properly, so this chapter should be read as a first, partial answer to that earlier call rather than a closed topic. The physical motivation running throughout, that type III factors are not a mathematical curiosity but the generic case for the local algebras of quantum field theory, where the modular flow of the vacuum state restricted to a wedge region is a Lorentz boost by the Bisognano-Wichmann theorem, is stated here only as motivation and is not developed, \cite[TODO: collocazione precisa]{TODO-BisognanoWichmann}, \cite[TODO: collocazione precisa]{TODO-Haag}.

What this chapter does establish is that the modular operator is not merely a technical device for proving the Tomita-Takesaki theorem, but the structure through which a von Neumann algebra, with no further input than a single state, already knows its own dynamics, its own equilibrium temperature, and its own fine classification. That three such different-looking facts come from one operator is the thread this chapter has tried to make visible.
\end{document}